\documentclass{report}
\usepackage[utf8]{inputenc}
\usepackage[italian]{babel}
\usepackage{amsmath, amssymb, amsfonts, mathtools}
\usepackage{float} % provides the H float specifier
\usepackage{booktabs} % provides \toprule, \midrule, \bottomrule used in tables
\usepackage{xcolor}
\usepackage{listings}
\usepackage{listings}
\usepackage{xcolor}
\usepackage[ruled,vlined]{algorithm2e}
\usepackage{hyperref}
\usepackage{csquotes}
\usepackage{comment}
\usepackage{amsthm}
\newtheorem{definition}{Definizione}[chapter] 
\usepackage{graphicx}
\usepackage{caption}
\captionsetup{hypcap=false}
%setkeys{Gin}{draft} % questo toglilo quando torni nell altro pc



\graphicspath{
    {/home/massi1/Documenti/uni/TESI/PARTE_SPERIMENTALE/results/detector/grafici/distanza_reali_random/}
    {/home/massi1/Documenti/uni/TESI/PARTE_SPERIMENTALE/results/generator/grafici/}
}

\setlength{\parindent}{0pt}
\lstset{
    basicstyle=\ttfamily\small,
    frame=single,
    breaklines=true,
    keywordstyle=\color{blue},
    commentstyle=\color{gray},
}


\begin{document}

\thispagestyle{empty}

\begin{center}

\vspace*{1cm}

{\Large \textsc{\mbox{Alma Mater Studiorum -- Università di Bologna}}}


\vspace{0.3cm}

\rule{\textwidth}{0.4pt}

\vspace{0.5cm}

{\large \textsc{Scuola di Scienze}}

\vspace{0.2cm}

{\large Corso di Laurea in Informatica}

\vspace{2.5cm}

\begin{center}
{\large\bfseries
Progettazione e Validazione di Tecniche\\[0.1cm]
per Generazione di File Casuali~Verosimili\\[0.1cm]
per Data Flooding against Ransomware
}
\end{center}






\vspace{3cm}

\begin{minipage}[t]{0.45\textwidth}
\raggedright
\textbf{Relatore:}\\
Prof.\\
SAVERIO GIALLORENZO\\[0.4cm]
\textbf{Correlatore:}\\
Dott.\\
SIMONE MELLONI
\end{minipage}
\hfill
\begin{minipage}[t]{0.45\textwidth}
\raggedleft
\textbf{Presentata da:}\\
MASSIMILIANO FIORAVANTI
\end{minipage}

\vspace{2.0cm}

{\large Sessione III}\\[0.2cm]
{\large Anno Accademico 2024/2025}

\end{center}

\clearpage


\clearpage

\begin{abstract}
Il ransomware rappresenta una delle minacce informatiche più diffuse e dannose
per utenti e organizzazioni, poiché la cifratura indiscriminata dei file può
causare gravi perdite economiche e operative. Tra le strategie di difesa proposte
in letteratura, il Data Flooding against Ransomware mira a ridurre l’impatto degli attacchi inducendo
il malware a cifrare file-esca.

Tra i metodi per la generazione dei file-esca ci sono approcci basati su file casuali.
Un problema di questa modalità è che file casuali che differiscono molto dai file reali
dell'utente possono essere discriminati e quindi ignorati dal ransomware, diminuendo
l'efficacia del contrasto.

\bigskip

Questo lavoro propone un sistema per la generazione di file-esca casuali,
progettati per risultare statisticamente compatibili con i file reali
dell’utente, senza ricorrere alla loro replica diretta. L’approccio si basa su una
modellazione statistica del contenuto binario dei file reali e sull’impiego di
metriche quantitative per valutare la verosimiglianza dei file generati rispetto a
specifici formati.

\bigskip

Viene inoltre sviluppato un sistema di valutazione automatico che consente di
stimare in modo oggettivo la verosimiglianza dei file generati e di validare
sperimentalmente l’efficacia dell’approccio proposto. I risultati ottenuti su
diversi formati di file mostrano che il metodo consente di generare file
artificiali difficilmente distinguibili da quelli reali mediante analisi
statistiche, rendendo l’approccio adatto al supporto di tecniche di difesa contro
attacchi ransomware.
\end{abstract}


\clearpage

\tableofcontents

\clearpage

\chapter{Introduzione}

\section{Ransomware}
Il termine Ransomware deriva dall'unione delle parole "ransom" (ostaggio) e "ware" (da software), indicando un programma che tiene in ostaggio dati o risorse digitali fino al pagamento di una somma di denaro.
In particolare, un ransomware è una tipologia di malware che, una volta infiltratosi in un sistema informatico, cifra i file della vittima, chiedendo successivamente un pagamento in cambio della chiave necessaria a decifrarli.
Quando un ransomware infetta un sistema, tende a colpire file sensibili come dati finanziari, documenti aziendali, database, file personali e altri contenuti di valore.

\bigskip
Negli ultimi anni il ransomware è diventato una delle minacce informatiche più diffuse e dannose a livello globale. Gli attacchi possono manifestarsi in vari modi, dai crypto-ransomware, i quali cifrano i file rendendoli inaccessibili, ai locker-ransomware, che bloccano l'accesso al sistema operativo stesso.
Questo tipo di malware rappresenta quindi un problema critico per utenti e organizzazioni, poichè può causare danni economici e operativi significativi.
La crescente professionalizzazione degli attaccanti e la diffusione del modello Ransomware-as-a-Service (RaaS) hanno ulteriormente amplificato la portata di questa minaccia \cite{alshaikh2020ransomware}.

\section{Contesto, Obiettivo e Contributo del Lavoro}

Il presente lavoro si colloca nell’ambito di \emph{Data Flooding against Ransomware} (DFaR),
un insieme di tecniche le quali prevedono la creazione, su richiesta,
di molti file-esca in cartelle ritenute sensibili, come ad esempio aree utente o, quando disponibile,
la posizione in cui il ransomware sta operando.
In questo modo la “trappola” viene portata vicino all’attaccante, aumentando la probabilità
di intercettare l’attività malevola e supportando non solo la rilevazione, ma anche la risposta
all’attacco.

\bigskip

Nella fase di mitigazione, il flooding mira a ostacolare un attacco in corso mescolando i file legittimi
con una grande quantità di esche, così che il malware sprechi tempo cifrando contenuti non rilevanti.

\bigskip

Tuttavia, la sola disponibilità di generare file-esca non è sufficiente: strategie basate su contenuti
puramente casuali possono produrre documenti facilmente distinguibili da file reali attraverso euristiche
e caratteristiche statistiche del contenuto. Ne consegue un problema per l’efficacia del data flooding:
i file generati devono risultare \emph{verosimili} rispetto a uno specifico formato almeno dal punto di vista
di proprietà osservabili a basso livello, così da aumentare la probabilità di non
essere scartati o ignorati da un attaccante. 

\bigskip

L'obiettivo del presente lavoro è lo studio e la progettazione di un sistema capace di generare file-esca casuali
verosimili, destinati a essere utilizzati in tecniche di data flooding contro attacchi ransomware.
L'obiettivo è produrre file che siano statisticamente simili a file reali,
in modo da aumentare la probabilità che vengano presi di mira dal ransomware nel suo processo di cifratura,
e diminuire la probabilità che il ransomware discrimini questi file e li ignori.

\bigskip

In dettaglio, il lavoro si propone di definire una nozione quantitativa di verosimiglianza dei file, 
da utilizzare come criterio di valutazione in grado di stabilire in modo automatico se 
un file sintetico possa essere considerato realistico o meno rispetto a uno specifico formato.

\bigskip

I contributi principali possono essere riassunti come segue:

\begin{itemize}
    \item la definizione di un insieme di metriche per il confronto statistico tra
    file reali e file generati, utilizzabili per stimare il grado di verosimiglianza di
    un file artificiale;
    \item la progettazione di un sistema di rilevamento che utilizza tali metriche
    come oracolo per la valutazione della verosimiglianza dei file;
    \item lo sviluppo di un metodo di generazione di file esca statisticamente
    compatibili con i vincoli imposti dal sistema di rilevamento;
    \item una valutazione sperimentale dell’approccio proposto su diversi
    formati di file e su dataset reali.
\end{itemize}


\section{Struttura della Tesi}

La tesi è strutturata come segue.

\bigskip
Nel Capitolo 2 vengono introdotti i concetti preliminari e il contesto teorico necessario 
alla comprensione del lavoro, con particolare riferimento alle principali tecniche di difesa 
contro il ransomware e all’approccio del data flooding.

\bigskip

Nel Capitolo 3 vengono presentate le rappresentazioni statistiche e le metriche utilizzate 
per descrivere e confrontare i file, insieme alle scelte adottate 
per selezionare le informazioni più significative da analizzare nei diversi formati di file.

\bigskip

Nel Capitolo 4 viene descritto il sistema di rilevamento (detector),
illustrando la definizione della verosimiglianza, il processo di ottimizzazione delle
soglie di rilevamento e i risultati sperimentali ottenuti.

\bigskip

Nel Capitolo 5 viene illustrato il sistema di generazione dei file verosimili,
basato su una modellazione statistica markoviana del contenuto binario 
e su un processo di generazione coerente con i vincoli imposti dal detector.

\bigskip

Infine, nel Capitolo 6 vengono discusse le conclusioni del lavoro, oltre a presentare possibili sviluppi futuri.




\chapter{Preliminari}

Questo capitolo introduce il contesto necessario a inquadrare il lavoro svolto.
In primo luogo viene fornita una panoramica delle principali tecniche di difesa e mitigazione
contro il ransomware, evidenziando obiettivi, punti di forza e limiti delle soluzioni più diffuse.
Successivamente l’attenzione si concentra sull’approccio del data flooding, e in particolare su
Data Flooding against Ransomware e sul tool Ranflood, che costituiscono il riferimento
applicativo di questa tesi. Il capitolo si chiude motivando il problema affrontato nel seguito:
la generazione di file-esca che non risultino facilmente distinguibili da quelli reali, introducendo
l’esigenza di criteri quantitativi per confrontare file reali e sintetici.


\section{Tecniche di Difesa contro il Ransomware}
La letteratura recente propone diversi metodi per rilevare e mitigare gli attacchi ransomware.
Per quanto riguarda la mitigazione, l'obiettivo fondamentale è quello di limitare i danni e recuperare efficacemente i dati compromessi.
Le tecniche più comuni includono:

\begin{itemize}
    \item \textbf{Segmentazione e isolamento della rete:} l’uso di Software Defined Networking (SDN) permette di bloccare dinamicamente i canali di comunicazione verso i server di comando e controllo (C2C), impedendo la cifratura remota dei dati \cite{cabaj2016sdn}.
    \item \textbf{Analisi e reverse engineering:} lo studio del comportamento del malware consente talvolta di recuperare i file cifrati, sfruttando debolezze nei meccanismi di cancellazione o generazione delle chiavi \cite{zimba2018data}.
    \item \textbf{Backup e ripristino automatico:} il mantenimento di copie di sicurezza periodiche e versioning consente di ripristinare lo stato precedente all’attacco, evitando il pagamento del riscatto \cite{baykara2018safezone}.
    \item \textbf{Sistemi di rilevamento basati su entropia o comportamento:} framework come quelli brevettati da McAfee e Dell monitorano la variazione dei file system (entropia, deduplicazione, snapshot) per rilevare cifrature anomale e ripristinare automaticamente le copie integre \cite{davies2022entropy}.
    \item \textbf{Data flooding:} l’iniezione di grandi quantità di dati innocui nei file critici può ostacolare l’efficacia della cifratura ransomware, rallentando il processo e aumentando la probabilità di rilevamento \cite{berardi2023dataflooding}.
\end{itemize}

In generale, si evidenzia che le tecniche di mitigazione sono efficaci solo se integrate in una strategia
più ampia che include prevenzione, aggiornamenti regolari e formazione degli utenti.
La difficoltà di sviluppare contromisure universali è legata alla
continua evoluzione dei ceppi e alle variabili umane che influenzano la risposta all’incidente \cite{kapoor2021review}.




\section{Data Flooding against Ransomware (DFaR) e Ranflood}
In questo lavoro, ci si concentra sull’approccio di mitigazione dell'attacco basato sul data flooding. 
L'idea alla base di questa tecnica è quella di inondare di file-esca il sistema non appena viene rilevata
un'attività sospetta riconducibile a un attacco ransomware. Così facendo, il ransomware viene indotto a
cifrare questi file, diluendo l'impatto dell'attacco sui dati reali e al contempo
rallentando il processo di cifratura.

\bigskip
Un tool che implementa tecniche di data flooding è Ranflood, il quale mette a disposizione diverse strategie, tra cui:
\begin{itemize}
  \item \textbf{Random:} genera file casuali nella posizione target per aumentare la superficie d'attacco.
  \item \textbf{On-the-fly:} replica file reali (pre-snapshot + firme) per creare copie valide.
  \item \textbf{Shadow:} utilizza snapshot/archivi preesistenti per mantenere copie di recovery.
\end{itemize} 

Concentrandoci sulla strategia Random, l'idea di generare file-esca casuali è concettualmente semplice e generalmente efficace. Tuttavia, alcuni lavori indicano una limitazione di questo approccio: Genç et al.~\cite{genc2019deception} mostrano come file-esca casuali possano essere individuati e ignorati da alcuni ransomware, ad esempio mediante controlli su metadati o statistiche di contenuto.
In particolare, file che presentano caratteristiche anomale, come un'entropia eccessivamente alta, possono essere facilmente identificati e esclusi durante il processo di cifratura.

\bigskip

Da quanto sopra, ne consegue la necessità di generare file credibili, esche che siano statisticamente simili ai file reali dell'utente, così da ridurre la probabilità che questi ultimi vengano ignorati.
Una via intermedia fra la generazione di file random e di file-copia di quelli reali è quindi quella di generare file simili ma non identici ai file reali, in modo da massimizzare la probabilità che il ransomware prenda di mira questi file.

\bigskip

In questo contesto, è quindi opportuno introdurre metriche che consentano di quantificare la similarità tra file reali e generati.
Nel presente lavoro si adotteranno rappresentazioni statistiche del contenuto binario (es. Byte Frequency Distribution) e misure di distanza/divergenza fra distribuzioni.
In particolare si farà riferimento a metriche già utilizzate in letteratura per il confronto tra file come entropia di Shannon, (per misurare il grado di casualità), la Jensen–Shannon Divergence (JSD) per la distanza simmetrica fra distribuzioni, la Manhattan distance, la Cosine Similarity e altre varianti (es. Total Variation Distance).
Tali metriche permetteranno di quantificare quanto un file generato artificialmente si allontani da un file reale.




\chapter{Metriche}

Questo capitolo introduce gli strumenti quantitativi utilizzati nel seguito per
descrivere un file in modo misurabile.

\bigskip

In primo luogo vengono presentate alcune rappresentazioni statistiche del contenuto binario
(BFD e N-gram), che permettono di ottenere un “profilo” numerico del file.

\bigskip

Successivamente si discute come rendere tale profilo più informativo nei formati che includono
porzioni compresse o poco significative, motivando le scelte adottate per concentrare l’analisi
sulle parti più rilevanti ai fini del confronto.

\bigskip

Infine vengono introdotte le metriche di distanza e divergenza impiegate per
quantificare la somiglianza tra profili statistici: queste misure costituiscono la base operativa
con cui verrà definita e valutata la verosimiglianza dei file nei capitoli successivi.

\section{Byte Frequency Distribution (BFD)}

Un file digitale può essere visto come una sequenza di byte, ciascuno dei quali assume un valore intero compreso tra 0 e 255.
La BFD rappresenta la distribuzione statistica di tali valori all’interno di un file, mostrando quante volte compare ciascun byte rispetto al totale dei byte contenuti.
In altre parole, si costruisce un array di 256 elementi (uno per ogni possibile valore di byte) e si conta il numero di occorrenze di ciascuno.

Formalmente, la BFD di un file $F$ è un vettore
\[
B = [f_0, f_1, \dots, f_{255}],
\]
in cui ogni elemento $f_i$ indica la frequenza relativa del valore di byte $i$:
\[
f_i = \frac{1}{N} \sum_{j=1}^{N} \delta(x_j = i)
\]
dove $x_j$ è il valore del byte in posizione $j$, $N$ è la lunghezza del file, e $\delta$ è la funzione di Kronecker.
\bigskip

Ogni tipo di file tende a produrre pattern caratteristici nella sua distribuzione globale dei byte, considerando l'intero contenuto del file.
Questi pattern possono essere considerati una firma statistica del formato.

McDaniel e Heydari~\cite{mcdaniel2003content} hanno mostrato come analizzando la distribuzione dei byte di più file appartenenti allo stesso formato,
si ottenga un profilo medio rappresentativo di quello specifico formato.
Gli autori affiancano a questo approccio anche un’analisi separata basata su header e trailer, trattata come algoritmo distinto.
Un file sconosciuto dunque può essere confrontato con tali profili per verificarne la similarità statistica.
Gli autori propongono inoltre di estendere questa rappresentazione, introducendo la Byte Frequency Cross-correlation (BFC),
la quale, oltre a contare la frequenza dei singoli byte, analizza le relazioni tra coppie di byte, consentendo di
cogliere pattern strutturali più complessi.


Il concetto di BFD può essere esteso anche oltre la classificazione per tipo, infatti
confrontando le distribuzioni di byte di due file è possibile stimare quanto i file siano statisticamente simili.
In tal senso, l’uso di metriche di distanza o divergenza consente di quantificare tale somiglianza, fornendo una base per il confronto fra file reali e generati.



\section{Rappresentazione N-gram}

Sebbene la Byte Frequency Distribution fornisca una rappresentazione utile del contenuto di un file, essa non cattura le relazioni sequenziali tra i byte,
in quanto considera solo la frequenza individuale di ciascun byte, ignorando le dipendenze locali tra byte consecutivi, 
i quali sono spesso rilevanti per descrivere la struttura interna del file.

Per superare questo problema, diversi lavori in letteratura propongono l’uso di rappresentazioni N-gram di byte,
nelle quali si considerano sequenze di $N$ byte consecutivi come unità di analisi \cite{li2019mdba,kolter2006executables}.

Come mostrano Zhang et al.~\cite{li2019mdba}, questo tipo di rappresentazione consente di catturare pattern più complessi e strutturati all’interno del file,
modellando pattern locali che riflettono meglio la natura del contenuto.

\bigskip

Formalmente, la rappresentazione N-gram di un file $F$ è ottenuta estraendo tutte le sequenze di $N$ byte consecutivi presenti nel file.
Ad ogni sequenza è associata una frequenza di occorrenza, ottenendo cosi un vettore
di lunghezza $256^N$ che descrive la distribuzione statistica delle sequenze di byte:
\[
G = [g_0, g_1, \dots, g_{256^N - 1}],
\]

\bigskip

Con particolare attenzione nel caso di $N=2$, si ottiene la rappresentazione a bigrammi, la quale cattura le relazioni tra coppie di byte consecutivi e 
quindi permette di modellare pattern più complessi rispetto alla semplice BFD.




\section{Estrazione Strutturale dei File (Structural Bytes)}

L'estrazione delle caratteristiche dai file può essere affrontata in diversi modi.
Il più semplice consiste nell'analizzare l'intero contenuto del file, calcolando la BFD o la N-gram su tutti i byte presenti.
Questo approccio, seppure semplice, non tiene conto della struttura interna del file, la quale può variare significativamente tra formati diversi.
Più in particolare, molti formati come per esempio PDF, JPG, DOCX presentano porzioni di dati compressi, non strutturati e dunque altamente entropici.
Queste sezioni tendono ad appiattire le distribuzioni dei byte, rendendo più difficile la distinzione tra file \cite{garfinkel2010future}.

\bigskip

Per arginare il problema, diversi lavori in letteratura come \cite{garfinkel2012dfxml} propongono di estrarre solo le porzioni 
strutturali del file, ovvero quelle sezioni che contengono metadati, header, marcatori e altre informazioni che riflettono la struttura del formato,
in modo tale da ottenere rappresentazioni statistiche più significative.

\bigskip

Considerando i formati analizzati in questo lavoro (PDF, JPG, DOCX, TXT), si è scelto di adottare un approccio simile,
estraendo le porzioni strutturali dei file per calcolare la N-gram.

Per esempio:
\begin{itemize}
    \item I PDF includono sezioni testuali, dizionari e metadati, mentre gli stream binari \texttt{(stream...endstream)} sono dati compressi senza una struttura particolarmente rilevante \cite{vanderknijff2009pdf}.
    \item Le immagini JPG contengono un header con marcatori rilevanti che definiscono segmenti strutturati del formato, mentre le sezioni dopo il segmento SOS (Start of Scan) sono altamente entropiche e dunque poco informative \cite{vandermeer2024jpeg}.
    \item Un file DOCX è caratterizzato da archivi ZIP i quali contengono sezioni XML strutturate, ed anche altre sezioni compresse ed altamente entropiche che possono essere ignorate \cite{garfinkel2009xml}.
\end{itemize}


\subsection{Estrazione Strutturale per PDF}

Per i PDF si rimuovono \texttt{stream...endstream}, in quanto contengono dati compressi e non strutturati (come per esempio immagini o font).

Dopo l'estrazione, si estraggono:
\begin{itemize}
    \item i 64KB iniziali del file contenente header e oggetti iniziali.
    \item gli 8KB finali del file, contenenti il trailer del documento.
\end{itemize}

Il processo descritto può essere sintetizzato come segue

\bigskip


\begin{algorithm}[H]
\caption{Estrazione strutturale PDF}

data $\gets$ read(path)\;
text $\gets$ decode(data, ``latin1'')\;
text $\gets$ remove\_streams(text)\;
head $\gets$ first $64$KB of text\;
tail $\gets$ last $8$KB of text\;
\Return encode(head + tail, ``latin1'')
\end{algorithm}


\subsection{Estrazione Strutturale per JPG}

Il formato JPG è caratterizzato da un header iniziale contenente marcatori strutturati (segmenti \texttt{FFxx}), seguito poi da un flusso compresso
di dati dopo il segmento \texttt{SOS} (Start of Scan).

Per ottenere una rappresentazione statistica significativa, si considerano:
\begin{itemize}
    \item l’header JPEG fino al segmento \texttt{SOS};
    \item una piccola finestra post–SOS contenente unicamente i marcatori \texttt{FFxx}.
\end{itemize}


L’algoritmo seguente sintetizza il procedimento adottato:

\bigskip

\begin{algorithm}[H]
\caption{Estrazione strutturale JPEG}

data $\gets$ read(path)\;

\If{data non inizia con SOI (\texttt{FF D8})}{
    \Return primi\ 256\ byte\;
}

sos $\gets$ indice del marker SOS (\texttt{FF DA})\;

\If{sos non trovato}{
    header $\gets$ data\;
}
\Else{
    header $\gets$ data fino a sos\;
}

markers $\gets$ estrai\_marcatori\_FFxx(data dopo sos, finestra = 1024)\;

\Return concatenate(header, markers)\;
\end{algorithm}



\subsection{Estrazione Strutturale per DOCX}

Questo formato è basato su archivi ZIP contenenti file XML strutturati e altre risorse. Per ottenere una rappresentazione significativa,
vengono selezionati solo i file XML rilevanti, quali:
\begin{itemize}
    \item \texttt{word/document.xml}
    \item \texttt{word/styles.xml}
    \item \texttt{file .rels}
    \item \texttt{[Content Types].xml}
\end{itemize}
vengono invece ignorati:
\begin{itemize}
    \item immagini \texttt{word/media/*}
    \item oggetti incorporati \texttt{embeddings/}
\end{itemize}

\begin{algorithm}[H]
\caption{Estrazione strutturale DOCX}

apri\_zip(path)\;
parts $\gets$ lista vuota\;

\ForEach{file all'interno dell'archivio}{
    nome $\gets$ lowercase(file)\;

    \If{nome inizia con \texttt{word/media/} \textbf{o} contiene \texttt{/media/} \textbf{o} \texttt{embeddings/}}{
        salta\;
    }

    \If{nome termina con \texttt{.xml} \textbf{o} \texttt{.rels} \textbf{o} è \texttt{[Content\_Types].xml}}{
        testo $\gets$ decode(file, UTF-8)\;
        aggiungi testo a parts\;
    }
}

text\_all $\gets$ concatenazione(parts)\;

\If{text\_all è vuoto}{
    \Return primi\ 4096\ byte\ del file\;
}

\Return encode(text\_all, ``latin1'')\;
\end{algorithm}



\subsection{Caso speciale: TXT}

I file TXT sono file di testo semplice, i quali non presentano sezioni compresse o non strutturate.
Per questo motivo non viene effettuata alcuna estrazione strutturale, in quanto l’intero contenuto del file è rilevante per la rappresentazione statistica.

\bigskip    

Può tuttavia succedere che il file considerato sia troppo piccolo (o addirittura vuoto).
In questi casi, è opportuno definire una soglia minima di lunghezza (256 byte), al di sotto della quale si raggiunge la lunghezza minima 
ripetendo il contenuto del file fino a raggiungere tale soglia.

\bigskip

Questo permette di ottenere una rappresentazione statistica più robusta, evitando che file troppo piccoli possano introdurre rumore o 
distorsioni nelle metriche calcolate in seguito.


\subsection{Approccio Generalizzato all’Estrazione Strutturale: Analisi e Limiti}

In una prima fase di analisi, si è valutata l’efficacia di un approccio generalizzato per l’estrazione strutturale dei file.
Più nel dettaglio, l'idea era quella di catturare i byte significativi del file in modo totalmente indipendente dal formato,
evitando così regole specifiche per ciascun tipo di documento.

Un primo approccio consiste nel combinare:
\begin{itemize}
    \item una porzione iniziale del file.
    \item una porzione finale.
    \item una selezione dei blocchi interni con entropia più bassa.
\end{itemize}

Lo scopo è quello di catturare le sezioni più strutturate del file, ignorando le parti compresse o altamente entropiche.

\bigskip

Tale approccio, sebbene interessante, dimostra diversi limiti rilevanti.
In primo luogo, essendo le sezioni strutturali fortemente dipendenti dal formato, l’estrazione generalizzata tende a perdere informazioni rilevanti,
portando a rappresentazioni statistiche meno significative.

\bigskip

In secondo luogo, la distinzione tra file reali e file random diventa meno marcata, riducendo l’efficacia delle metriche di confronto.

\bigskip

Infine, per rendere efficace una selezione automatica delle regioni ad alta informazione è necessario disporre di dataset molto ampi per ciascun formato,
condizione non soddisfatta nel contesto di questo lavoro, nel quale sono stati utilizzati dataset che contengono 
approssimativamente dai 10\,000 ai 5\,000 file per formato.

\bigskip

Per queste ragioni, si è deciso di adottare l’approccio basato su regole specifiche per ciascun formato,
come descritto in precedenza.


\section{Metriche di Distanza e Divergenza}

Una volta ottenuta la rappresentazione statistica del file sotto forma di BFD (o BFC), è necessario disporre di metriche quantitative per confrontarle.
Lavori in letteratura (come per esempio Tabish et al.~\cite{tabish2009malware}) utilizzano metriche di distanza e divergenza per quantificare quanto due distribuzioni di byte differiscono tra loro.
Queste misure consentono di stimare la similarità statistica tra file, ovvero quanto due file presentano un contenuto binario strutturalmente simile o dissimile.

\subsection{Entropia di Shannon}

Una prima metrica utile per descrivere o classificare file è l’entropia di Shannon, la quale misura il grado di casualità della distribuzione dei byte.
In particolare, valori elevati indicano un contenuto casuale, al contrario, valori bassi riflettono un contenuto più strutturato.

Sorokin~\cite{sorokin2011entropy} introduce il concetto di \emph{structural entropy} per il confronto di file in base alla distribuzione locale dell’entropia, mostrando come la sequenza dei valori entropici può essere vista come una firma statistica della sua struttura interna.

Formalmente, l’entropia di Shannon di una distribuzione di byte è definita come:
\[
H = - \sum_{i=0}^{255} p(i) \log_2 p(i)
\]
dove $p(i)$ rappresenta la frequenza relativa del byte $i$ all’interno della finestra di analisi.

\subsection{Jensen–Shannon Divergence (JSD)}

Un’altra metrica utilizzata per tale scopo è la \emph{Jensen–Shannon Divergence} (JSD).
A differenza della KL divergence, la JSD è simmetrica e sempre definita, rendendola stabile.

La JSD viene impiegata per quantificare quanto due distribuzioni di probabilità differiscono o si assomigliano, e quindi può essere utilizzata per misurare quanto le BFD di due file siano statisticamente simili.

Formalmente, la JSD è definita come:
\[
\mathrm{JSD}(P,Q) = \frac{1}{2}\mathrm{KL}(P\|M) + \frac{1}{2}\mathrm{KL}(Q\|M)
\]
dove $M = \frac{1}{2}(P+Q)$, $\mathrm{KL}$ è la divergenza di Kullback–Leibler e $P$, $Q$ rappresentano le due distribuzioni di byte da confrontare.

Più nel dettaglio, il valore di JSD è compreso fra 0 e 1, dove valori vicini a 0 indicano distribuzioni molto simili, al contrario valori vicini a 1 riflettono distribuzioni diverse.

Tabish et al.~\cite{tabish2009malware} utilizzano la JSD per analizzare la differenza informativa tra distribuzioni di byte in file benigni e file malevoli, mostrando che la divergenza cresce proporzionalmente alla differenza di struttura statistica del contenuto.


\subsection{Manhattan Distance}

La Manhattan distance è una misura di dissimilarità fra due distribuzioni, calcolata come somma delle differenze assolute tra le rispettive sequenze di byte:
\[
D_{\mathrm{L1}}(X,Y) = \sum_{i=0}^{255} |X_i - Y_i|
\]
Nella letteratura, è usata in diversi lavori che si occupano di file type classification e content-based analysis.
In particolare, Amirani et al.~\cite{amirani2008content} la utilizzano per confrontare le distribuzioni di byte di file e frammenti, mostrando che valori bassi di Manhattan distance riflettono un’elevata similarità fra i file.

\subsection{Cosine Similarity}

Un’ulteriore metrica utile è la \emph{Cosine Similarity}, una misura di tipo angolare la quale valuta quanto due distribuzioni siano orientate nella stessa direzione.
Riguardo alla BFD, permette di misurare quanto due file presentino pattern di frequenze di byte simili.

La Cosine Similarity fra due distribuzioni $X$ e $Y$ è definita come:
\[
\mathrm{cos}(X,Y) = \frac{\sum_{i=0}^{255} X_i Y_i}{\sqrt{\sum_{i=0}^{255} X_i^2}\sqrt{\sum_{i=0}^{255} Y_i^2}}
\]
Il risultato varia fra 0 e 1, in particolare valori vicini ad 1 indicano distribuzioni quasi identiche, mentre valori vicini allo 0 distribuzioni dissimili.

In letteratura, Ahmed et al.~\cite{ahmed2010cosine} utilizzano questa metrica nel contesto del \emph{content-based file type identification}, confrontando gli istogrammi di frequenza dei byte per classificare i file in base al loro tipo.
Anche Tabish et al.~\cite{tabish2009malware} includono la Cosine Similarity per misurare la differenza tra distribuzioni di byte, evidenziando come la correlazione angolare sia un indicatore efficace per distinguere file appartenenti a categorie diverse.


\subsection{Total Variation Distance (TVD)}

La Total Variation Distance (TVD) è una metrica probabilistica che quantifica la massima differenza tra due distribuzioni di probabilità.
Nel contesto delle BFD, essa misura quanto la distribuzione dei byte di un file differisce da quella di un altro file, fornendo una stima diretta del grado di dissimilarità statistica tra i due contenuti.\\
Formalmente è definita come:
\[
\mathrm{TVD}(P,Q) = \frac{1}{2} \sum_{i=0}^{255} |P_i - Q_i|
\]     
dove $P$ e $Q$ sono le distribuzioni di probabilità dei byte dei due file.
Il valore della distanza varia tra 0 e 1, con 0 che indica distribuzioni identiche e 1 che rappresenta distribuzioni completamente disgiunte.



\chapter{Sistema di Rilevamento dei File Verosimili (Detector)}

Nella sezione precedente sono state presentate diverse metriche utili a descrivere e confrontare file in base alla loro distribuzione di byte.

\`{E} quindi necessario introdurre un sistema di rilevamento con l'obiettivo di valutare la verosimiglianza di file generati artificialmente.
A tale scopo si utilizzeranno le metriche descritte per confrontare i file generati con delle soglie predefinite, stabilite analizzando un dataset di file reali.

Di seguito verranno descritte le varie fasi del detector, descrivendo le scelte implementative adottate e gli strumenti utilizzati.


\section{Analisi Preliminare: Distanza tra File Reali e File Random}

Come primo passo per la costruzione del detector, è stata condotta un'analisi per valutare la distanza statistica fra i file reali e file generati casualmente.
Lo scopo di questa prima analisi è quello di quantificare la differenza fra i file reali e random, facendo vedere che, a livello statistico, i file random si discostano significativamente dai file reali.

Per farlo, si confrontano le rappresentazioni N-gram dei file reali con quelle dei file random mediante le metriche discusse in precedenza, ottenendo una stima quantitativa della distanza media tra i due insiemi.
\bigskip

In particolare, per questa fase è stato sviluppato uno script che, dati i dataset di file reali e di file random per ogni formato considerato,
produce un insieme di confronti appaiati tra i due dataset.

\bigskip

Per ogni formato considerato, lo script esegue le seguenti operazioni:
\begin{itemize}
    \item seleziona casualmente una coppia di file dello stesso formato (reale, randomico) dai dataset opportuni;
    \item calcola la N-gram dei due file selezionati;
    \item applica le metriche per confrontare le N-gram calcolate in precedenza;
    \item ripete il processo per un numero prefissato di iterazioni (20 000 nel nostro caso), per ottenere una stima statistica robusta.
\end{itemize}

Questa analisi evidenzia una chiara separazione statistica tra i file reali e quelli generati casualmente,
che mostrano distribuzioni di byte nettamente differenti.
\bigskip


La Figura \ref{fig:jsd_distance} mostra la statistica ottenuta con la metrica JSD calcolata sulle coppie file reale--file random per ciascun formato.
Similmente, la Figura \ref{fig:l1_distance} mostra i risultati ottenuti con la Manhattan distance,
la Figura \ref{fig:tvd_distance} con la Total Variation Distance,
e la Figura \ref{fig:cosine_similarity} con la Cosine Similarity.

\bigskip

In tutti i casi, si osservano valori di distanza/divergenza elevati, indicando che i file random si discostano significativamente dai file reali.

\bigskip

La Figura \ref{fig:pairwise_entropy} mostra invece i valori medi di entropia dei file reali e random.
Anche qui emerge una differenza significativa, con i file randomici che presentano un'elevata entropia rispetto ai file reali,
i quali mostrano valori più bassi, a riflettere una struttura più ordinata e definita.


\begin{figure}[htbp]
    \centering
    \includegraphics[width=0.8\textwidth]{jsd_distance_stats_gray.pdf}
    \caption{Distanza media tra file reali e file random calcolata mediante la metrica JSD.}
    \label{fig:jsd_distance}
\end{figure}


\begin{figure}[htbp]
    \centering
    \includegraphics[width=0.8\textwidth]{l1_distance_stats_gray.pdf}
    \caption{Distanza media tra file reali e file random calcolata mediante la Manhattan Distance.}
    \label{fig:l1_distance}
\end{figure}


\begin{figure}[htbp]
    \centering
    \includegraphics[width=0.8\textwidth]{tvd_distance_stats_gray.pdf}
    \caption{Distanza media tra file reali e file random calcolata mediante la Total Variation Distance (TVD).}
    \label{fig:tvd_distance}
\end{figure}


\begin{figure}[htbp]
    \centering
    \includegraphics[width=0.8\textwidth]{cosine_distance_stats_gray.pdf}
    \caption{Distanza media tra file reali e file random calcolata mediante la Cosine Similarity.}
    \label{fig:cosine_similarity}
\end{figure}


\begin{figure}[htbp]
    \centering
    \includegraphics[width=0.8\textwidth]{entropy_real_vs_random.pdf}
    \caption{Entropia media dei file reali e dei file random.}
    \label{fig:pairwise_entropy}
\end{figure}

\bigskip

Poiché i valori osservati negli istogrammi sono concentrati in prossimità
del massimo, le differenze tra i formati risultano poco evidenti a livello visivo.
Per rendere il confronto più chiaro, nella Tabella~\ref{tab:mean-real-random} sono riportati i valori medi delle metriche per ciascun formato.





\begin{table}[H]
\centering
\caption{Statistiche medie per formato (real-random)}
\label{tab:mean-real-random}
\renewcommand{\arraystretch}{1.25}
\setlength{\tabcolsep}{8pt}
\begin{tabular}{lcccccc}
\toprule
\textbf{Formato} & \textbf{JSD} & \textbf{TVD} & \textbf{L1} & \textbf{Cosine\_sim} & \textbf{Entropy\_random} & \textbf{Entropy\_real} \\
\midrule
docx & 0.9792 & 0.9882 & 1.9765 & 0.0095 & 11.9379 & 8.1562 \\
jpg  & 0.9833 & 0.9948 & 1.9895 & 0.0139 & 14.4502 & 6.7402 \\
pdf  & 0.9710 & 0.9893 & 1.9785 & 0.0237 & 15.2401 & 7.4579 \\
txt  & 0.9745 & 0.9911 & 1.9821 & 0.0289 & 15.5480 & 7.1467 \\
\bottomrule
\end{tabular}
\end{table}



\section{Definizione della Verosimiglianza dei File}

Visti i risultati ottenuti nella fase preliminare, che mostrano come i file reali differiscano in modo significativo da quelli generati casualmente,
la domanda successiva da porsi è quando un file sintetico possa essere considerato "realistico" o "verosimile".
In altre parole, quanto un file generato artificialmente debba essere simile a un file reale per essere considerato credibile.

\bigskip

Per rispondere alla domanda, è necessario definire delle soglie di verosimiglianza per ciascuna metrica considerata,
al di sotto (o al di sopra, a seconda della metrica che si considera) delle quali un file generato viene considerato realistico.

\bigskip

Nelle sezioni seguenti viene illustrato l’approccio adottato per determinare tali soglie e costruire un sistema di decisione robusto.


\section{Costruzione della Firma Media}

Per poter stabilire se un file generato sia o meno verosimile, è necessario disporre di un riferimento statistico rappresentativo dei file reali.
A tale scopo viene costruita una \emph{firma media} per ciascun formato considerato, la quale riassume la struttura tipica dei file reali.

\bigskip

La firma media viene calcolata come la media aritmetica delle N-gram dei file reali presenti nel dataset.
Il vettore risultante corrisponde dunque alla distribuzione media dei byte per quel formato specifico.

Tale firma media funge da riferimento per confrontare la verosimiglianza di un file: un file che presenta una distribuzione simile alla firma media viene considerato realistico.


\bigskip

Una volta definita la firma media, per ciascun file è possibile calcolare la distanza fra la sua distribuzione e la firma media,
ottenendo un insieme di valori per ogni metrica considerata.

Resta quindi da stabilire quali valori di distanza/divergenza possono essere considerati accettabili per definire un file come verosimile.

\bigskip

Nella sezione che segue viene descritto l’approccio adottato per determinare tali soglie di verosimiglianza.




\section{Definizione e Ottimizzazione delle Soglie di Verosimiglianza}

Come descritto sopra, l'obiettivo è stabilire, per ogni metrica, una \emph{soglia di verosimiglianza} tale per cui 
un file è considerato verosimile se la distanza dalla sua firma media è inferiore (o superiore) a tale soglia.

\bigskip

Un primo metodo per determinare queste soglie è un approccio empirico, basato sull'analisi statistica dei valori di distanza calcolati sui file reali.
Ad esempio, si possono calcolare la media e la deviazione standard delle distanze dei file reali dalla firma media, e definire la soglia come:
\[\text{soglia} = \mu + k \cdot \sigma
\]
dove $k$ è un fattore di scala scelto in base al livello di tolleranza desiderato.

Tuttavia, questo approccio presenta alcuni limiti, ad esempio la scelta di $k$ può essere arbitraria e non garantisce una separazione ottimale tra file reali e generati.

\bigskip

Pertanto, per individuare soglie più robuste, è necessario adottare un approccio di ottimizzazione più sistematico.

\bigskip

Per stabilire una soglia ottimale, si rende necessario disporre di un indicatore quantitativo che misuri la qualità della classificazione ottenuta.

Nel nostro caso, fissata una soglia $T$, possiamo classificare un file come realistico o non realistico a seconda della distanza 
dalla firma media.

\bigskip

Questo porta alla definizione di quattro insiemi:
\begin{itemize}
    \item \(TP\): l’insieme dei file reali classificati come realistici (Veri Positivi);
    \item \(TN\): l’insieme dei file randomici classificati come non realistici (Veri Negativi);
    \item \(FN\): l’insieme dei file reali classificati come non realistici (Falsi Negativi);
    \item \(FP\): l’insieme dei file randomici classificati come realistici (Falsi Positivi).
\end{itemize}

\bigskip

L'idea è quindi quella di costruire una funzione obiettivo, che, dato uno specifico valore di soglia \(T\), misuri la qualità della classificazione ottenuta.

\bigskip

Dobbiamo dunque cercare di classificare tutti i file reali come realistici (massimizzare \(TP\)) ed in contemporanea evitare di accettare come realistici file randomici (minimizzare \(FP\)).
Pertanto, una possibile funzione obiettivo può essere definita come:
\[
C(T) = \alpha \cdot FN(T) + \beta \cdot FP(T)
\]

dove \(\alpha\) e \(\beta\) sono pesi che bilanciano l'importanza di minimizzare i falsi negativi e i falsi positivi,
e \(FN(T)\) e \(FP(T)\) rappresentano rispettivamente la percentuale di falsi negativi e falsi positivi ottenuti con la soglia \(T\), 
dividendo il numero di questi ultimi per il numero totale di file reali e randomici rispettivamente.

\bigskip

Una volta definita la funzione obiettivo \(C(T)\), il problema si riduce a trovare la soglia \(T^*\) per cui \(C(T)\) è minima.
Nel presente lavoro, per risolvere questo problema di ottimizzazione, si è scelto di utilizzare l'algoritmo di \emph{Dual Annealing}, descritto in seguito.


\subsection{Ricerca della Soglia Ottimale tramite Dual Annealing}

Per individuare la soglia ottimale \(T^*\) che minimizza \(C(T)\), può risultare inefficace utilizzare metodi di ottimizzazione classici come ad esempio la discesa del gradiente,
in quanto la funzione obiettivo può avere molteplici minimi locali, derivanti dalla distribuzione non lineare degli score di distanza.

\bigskip

Per questa ragione si è scelto di adottare un approccio di ottimizzazione globale basato sull'algoritmo di \emph{Dual Annealing}.

\bigskip

Quest'algoritmo combina elementi del Classical Simulated Annealing e del Fast Simulated Annealing in un’unica procedura di ottimizzazione globale.
Come riportato nella documentazione ufficiale di SciPy \cite{scipy_dual_annealing}, Dual Annealing usa una particolare visiting distribution di tipo Cauchy-Lorentz modificata,
che permette di effettuare salti esplorativi anche molto ampi nello spazio delle soluzioni, in modo da minimizzare
la probabilità di finire in eventuali minimi locali.

\bigskip

In particolare, ad ogni iterazione il metodo genera una nuova soluzione candidata \(T\) applicando la visiting distribution e decide 
se accettarlo attraverso una probabilità di accettazione che dipende dalla differenza di costo e dalla temperatura artificiale.

Grazie a queste tecniche, l'algoritmo di Dual Annealing è in grado di esplorare efficacemente lo spazio delle soluzioni,
consentendo di individuare la soglia ottimale \(T^*\).


\section{Overfitting e Cross-validation}

Arrivati a questo punto, è importante considerare alcuni aspetti critici che possono influenzare l’affidabilità e la robustezza 
del modello sviluppato. In particolare, il processo di ottimizzazione delle soglie può
essere soggetto a fenomeni di overfitting quando il dataset disponibile è limitato.

\bigskip

Nel contesto di questo lavoro, si dispone infatti di un numero relativamente basso di file per ciascun formato.
Questo può portare a situazioni in cui si hanno soglie di verosimiglianza che si adattano in modo troppo specifico ai file presenti nel dataset,
riducendo la capacità del modello di generalizzare a file diversi.

\bigskip

Per ridurre il rischio di overfitting, è possibile adottare strategie che permettano di costruire dati su porzioni diverse dei dataset
disponibili, in modo da evitare che i risultati ottenuti siano troppo specifici a una determinata selezione di file.

\bigskip

Per questo motivo, nel presente lavoro è stato utilizzato un approccio di \emph{k-fold cross-validation}.

\bigskip

La cross-validation è una tecnica statistica di valutazione del modello. Essa suddivide l'intero dataset
in \(k\) sottoinsiemi (fold) di dimensioni uguali. Il modello viene poi addestrato su \(k-1\) fold e testato sul fold rimanente, 
ripetendo il processo \(k\) volte, in modo che ogni fold venga utilizzato una volta come set di test.

I risultati ottenuti vengono poi mediati per ottenere una stima più robusta delle prestazioni del modello.

Come descritto da Arlot e Celisse~\cite{arlot2010crossvalidation}, questo approccio permette di fornire una stima più affidabile della capacità di generalizzazione del modello,
riducendo il rischio di overfitting dovuto alla limitata dimensione del dataset.




\section{Implementazione del Sistema di Rilevamento}

Concluse le considerazioni metodologiche, nella seguente sezione sarà descritto nel dettaglio il sistema di rilevamento implementato,
mostrando le scelte tecniche adottate, il flusso di lavoro e gli strumenti utilizzati.

\bigskip

Il sistema di rilevamento è progettato con lo scopo di analizzare diversi tipi di file (PDF, JPG, DOCX, TXT) 
e valutare la loro verosimiglianza in base alle metriche e soglie definite in precedenza.

Il sistema si articola in quattro fasi principali:
\begin{itemize}
    \item Estrazione strutturale del file.
    \item Calcolo della firma media e degli score di distanza.
    \item Ottimizzazione delle soglie di verosimiglianza.
    \item Valutazione delle performance del detector.
\end{itemize}

Nelle sezioni seguenti verranno descritte in dettaglio ciascuna di queste fasi.





\subsection{Struttura Generale}

Il sistema di rilevamento è stato organizzato come una serie di moduli, ciascuno responsabile di una specifica fase del processo.

La struttura generale del sistema può essere sintetizzata come segue:

\begin{enumerate}

\item {Generazione dello split \(k\)-fold}.

Per ciascun formato si
costruiscono due insiemi separati: file reali e file randomici.  
Ogni insieme viene mescolato in modo casuale ma riproducibile (seed fisso) e poi suddiviso
in \(k\) blocchi disgiunti di dimensione il più possibile uniforme (la differenza tra blocchi è al massimo 1 file).

Nel fold \(i\)-esimo, il blocco \(i\) (sia reali sia randomici) costituisce il test, mentre l’unione dei restanti \(k-1\) blocchi costituisce il train.

    \item {Calcolo firma media e score di distanza per ciascun fold}.

Per ogni fold, i file che appartengono al fold non utilizzato per il test vengono usati per l'addestramento.
Per fare un esempio, nel caso in cui \(k=3\), si usano i file dei fold 1 e 2 per l'addestramento e i file del fold 3 per il test,
e così via per gli altri fold.

Su questo insieme di addestramento, viene calcolata la firma media per ciascun formato.
Successivamente, per ogni file dei fold di train e test, vengono calcolate le metriche di distanza rispetto alla firma media.

    \item {Ottimizzazione delle soglie di verosimiglianza per ciascun fold}.

Usando come base gli score ottenuti nel passo precedente sui file di addestramento, viene eseguita 
l'ottimizzazione delle soglie per ogni metrica e formato.

Come descritto in precedenza, l'ottimizzazione viene formulata come un problema di minimizzazione di una funzione obiettivo \(C(T)\),
al fine di bilanciare i falsi positivi e negativi.

    \item {Valutazione delle performance del detector}.

In questa fase le soglie ottimizzate con il fold di addestramento vengono applicate ai file del fold di test.

In particolare, per ciascun formato vengono calcolate:
\begin{itemize}
        \item il numero di falsi positivi (FP): file randomici classificati come realistici.
        \item il numero di falsi negativi (FN): file reali classificati come non realistici.
        \item il contributo di ciascuna metrica al processo decisionale.
\end{itemize}

I risultati dei singoli fold vengono aggregati per ottenere una stima complessiva delle performance del detector.

\end{enumerate}





\subsection{Calcolo della Firma Media e degli Score di Distanza}

Una volta definita la suddivisione in fold, il passo successivo consiste nel calcolare la firma media per ciascun formato
e gli score di distanza per ciascun file.

\bigskip

Il primo passo consiste nel convertire ciascun file nella sua rappresentazione numerica basata sulle N-gram, come descritto nella Sezione 3.2.
Nel presente lavoro è stato scelto di utilizzare $n=2$ (bigrammi), ottenendo un vettore di dimensione 65536 ($256^2$).

\bigskip

Con il seguente pseudocodice viene illustrato il processo di calcolo della N-gram per un file:


\begin{algorithm}[H]
\caption{BFD-2gram}
\SetKwInOut{Input}{Input}
\Input{contenuto binario del file ($b$)}

\If{$b$ è vuota}{
    \Return vettore di 65\,536 zeri\;
}

arr $\gets$ converti $b$ in array di interi 8-bit\;
\If{lunghezza(arr) $<$ 2}{
    \Return vettore di 65\,536 zeri\;
}

\For{$i = 0$ \KwTo lunghezza(arr) - 2}{
    idx $\gets$ arr[i] * 256 + arr[i+1]\;
    BFD[idx] $\gets$ BFD[idx] + 1\;
}

normalizza(BFD)\;
\Return BFD\;
\end{algorithm}


\bigskip

Dopo aver calcolato la N-gram per ciascun file, si procede al calcolo della firma media, la quale è ottenuta
come la media aritmetica delle N-gram dei file reali presenti nel fold di addestramento.

\bigskip

\begin{algorithm}[H]
\caption{CalcoloCentroide}
\KwIn{Lista di file $F_{\mathrm{real}}$}

\ForEach{$f \in F_{\mathrm{real}}$}{
    v $\gets$ BFD-2gram(f)\;
    lista\_vettori.append(v)\;
}

centroide $\gets$ mean(lista\_vettori)\;
normalizza(centroide)\;
\Return centroide\;
\end{algorithm}


\bigskip

Si procede poi al calcolo degli score di distanza per ciascun file, confrontando la sua N-gram con la firma media calcolata in precedenza.

\bigskip

\begin{algorithm}[H]
\caption{CalcoloScore}
\KwIn{Insieme di file $F = F_{\mathrm{real}} \cup F_{\mathrm{rand}}$, centroide $C$}

\ForEach{file $f \in F$}{
    v $\gets$ BFD\_2gram(f)\;

    score[f].jsd $\gets$ JensenShannon$(v, C)$\;
    score[f].tvd $\gets$ TotalVariation$(v, C)$\;
    score[f].l1  $\gets \|v - C\|_1$\;
    score[f].cos $\gets$ CosineSimilarity$(v, C)$\;
    score[f].ent $\gets$ Entropy$(v)$\;
}

\Return score\;
\end{algorithm}

\bigskip

Per ciascun fold, vengono prodotti due file csv:
\begin{itemize}
    \item \texttt{train\_scores\_fold\_i.csv}: contiene gli score di distanza per i file dei fold di addestramento.
    \item \texttt{test\_scores\_fold\_i.csv}:  contiene gli score di distanza per i file del fold di test.
\end{itemize}

Gli score calcolati in questa fase verranno poi utilizzati per l'ottimizzazione delle soglie di verosimiglianza e la valutazione delle performance del detector nelle fasi successive.



\subsection{Applicazione dell’Ottimizzazione delle Soglie}

Una volta ottenuti gli score di distanza per i file di addestramento, la fase successiva consiste
nell'individuare, per ciascuna metrica e ciascun formato, la soglia di verosimiglianza ottimale.

L’ottimizzazione viene eseguita in modo indipendente per ciascun fold,
considerando separatamente ogni combinazione formato–metrica.
Per ogni metrica, gli score dei file reali e randomici appartenenti al training set
vengono utilizzati per stimare una soglia che massimizzi la capacità discriminante
del modello.

\bigskip

La funzione obiettivo \(C(T)\) è la stessa funzione definita nella Sezione 4.4, ed è basata su una combinazione 
pesata di falsi positivi e falsi negativi. Nell'implementazione, questi valori sono calcolati 
a partire dagli score di distanza ottenuti nella fase precedente.

\bigskip

Per ciascuna metrica, il calcolo dei falsi positivi e negativi avviene applicando regole di decisione 
coerenti con la natura della metrica considerata. Ad esempio, per metriche di distanza come JSD, TVD e L1,
un file viene classificato come realistico se il suo score è inferiore alla soglia \(T\).

\bigskip

Come descritto in precedenza, per risolvere il problema di ottimizzazione si è scelto di utilizzare l'algoritmo di Dual Annealing,
il quale esplora lo spazio delle possibili soglie \(T\) e restituisce quella che minimizza la funzione obiettivo \(C(T)\).

\bigskip

Il risultato di questa fase è un insieme di soglie ottimali per ciascun formato e metrica,
che verranno poi utilizzate nella fase di valutazione delle performance del detector.




\subsection{Valutazione delle Performance}

Una volta ottenute le soglie ottimizzate, il passo finale per la costruzione del detector consiste nella valutazione delle sue performance
sui file del fold di test.

In questa fase, le soglie ottenute sul training set vengono applicate esclusivamente ai file del test set.

\bigskip

Per ognuno dei file di test, si caricano gli score di distanza calcolati e si confrontano con le soglie ottenute.

La decisione finale sulla verosimiglianza è ottenuta combinando i risultati delle singole metriche.
In particolare, la decisione è ottenuta applicando \emph{AND logico}: un file è considerato verosimile solamente se
tutte le metriche considerate rientrano nei rispettivi parametri di soglia.

Nel caso in cui anche una sola metrica non rientri nella soglia, il file viene classificato come non realistico.

\bigskip

La scelta di adottare un AND logico è motivata dalla necessità di ridurre al minimo la probabilità di falsi positivi.
Nel contesto di questo lavoro, infatti, l’accettazione di file non realistici potrebbe
compromettere l’efficacia del sistema, poiché tali file rischierebbero di essere facilmente individuati
e ignorati dal malware.

\bigskip

Questo approccio privilegia quindi una classificazione restrittiva, accettando un lieve incremento dei falsi negativi
come compromesso coerente con lo scenario applicativo considerato, in cui è preferibile scartare alcuni file
potenzialmente realistici piuttosto che accettare file statisticamente anomali.


\bigskip

Per quanto riguarda le singole metriche, come accennato prima, la regola di decisione dipende dalla natura della metrica stessa.

Nel dettaglio:
\begin{itemize}

    \item Per metriche di distanza/divergenza (JSD, TVD, $L_1$):

Un file è considerato realistico se il suo score è \emph{minore o uguale} alla soglia di
verosimiglianza stimata per il formato considerato;

Per esempio, se la soglia di verosimiglianza per JSD calcolata su un file PDF è 0.95, e la soglia calcolata per formato PDF e metrica JSD è 0.92,
allora il file viene considerato non realistico rispetto alla metrica JSD.

    \item Per la metrica di Cosine Similarity:

Un file è considerato realistico se il suo score è \emph{maggiore o uguale} alla soglia di
verosimiglianza stimata per il formato considerato;

    \item Per la metrica di Entropia:

la verosimiglianza viene valutata confrontando l'entropia del file con un intervallo accettabile $[T_{\text{low}}, T_{\text{high}}]$,
all'interno del quale un file realistico dovrebbe cadere.



\end{itemize}



\subsection{Detector Finale}

Una volta descritte le fasi di addestramento, ottimizzazione e valutazione, il sistema è utilizzabile 
in modalità operativa per valutare la verosimiglianza di nuovi file generati.

Le soglie ottenute dai vari fold vengono aggregate al fine di ottenere un set finale di soglie da utilizzare nel detector finale.

In particolare, per le metriche basate su una singola soglia (JSD, TVD, $L_1$, Cosine Similarity),
la soglia finale è ottenuta come il 95-esimo percentile delle soglie ottenute sui singoli fold.

Per la metrica di entropia, invece, gli intervalli finali $[T_{\text{low}}, T_{\text{high}}]$ sono ottenuti
calcolando il valore medio dei diversi fold.

\bigskip

Il detector finale opera come segue:

\bigskip

\begin{algorithm}[H]
\caption{DetectorOperativo}
\SetKwInOut{Input}{Input}

\Input{file $f$, centroidi $\{C_{fmt}\}$, soglie $\{T_{fmt}\}$}

fmt $\gets$ determinaFormato$(f)$\;

\If{fmt non supportato}{
    \Return NON\_PLAUSIBILE\;
}

$v \gets$ BFD-2gram$(f)$\;
$C \gets C_{fmt}$\;
$T \gets T_{fmt}$\;

\ForEach{metrica $m$}{
    $s_m \gets$ calcolaScore$(v, C, m)$\;
    verifica $s_m$ rispetto alla soglia $T_m$\;
}

\If{tutte le metriche soddisfano le soglie}{
    \Return PLAUSIBILE\;
}
\Else{
    \Return NON\_PLAUSIBILE\;
}
\end{algorithm}






\section{Valutazione del Detector}

In questa sezione si  presentano e analizzano i valori sperimentali ottenuti.

L'obiettivo di questa valutazione è misurare la capacità del modello di classificare file come realistici o randomici,
analizzando le prestazioni ottenute in termini di falsi positivi, falsi negativi e al contributo delle singole metriche.

\bigskip

La valutazione è stata condotta utilizzando una procedura di cross-validation
k-fold con $k = 3$, applicando a ciascun fold le soglie ottimizzate sul
corrispondente insieme di addestramento.


\subsection{Setup Sperimentale}

La valutazione utilizza più dataset, divisi per formato di file (PDF, JPG, DOCX, TXT).

In particolare, per ciascun formato si hanno due dataset distinti, uno contenente file reali e l'altro contenente file generati casualmente.

I file reali rappresentano documenti legittimi provenienti da dataset pubblici ampiamente utilizzati in ambito forense e di analisi dei file. In particolare, per i formati PDF e DOCX sono stati utilizzati documenti provenienti dai dataset GovDocs1~\cite{garfinkel2009govdocs} e NapierOne~\cite{davies2022napierone,napierone_dataset}.  
Per il formato JPG sono state utilizzate immagini provenienti da dataset pubblici di immagini naturali, mentre per i file TXT sono stati utilizzati documenti testuali provenienti da collezioni pubbliche di testi, come il dataset Project Gutenberg~\cite{bejan2020gutenberg}.

I file randomici sono stati invece generati utilizzando il tool Ranflood, il quale si basa sulla tecnica del Xorshift.

Per i formati PDF, JPG e TXT, il dataset è composto da circa 10\,000 file reali e
2\,000 file randomici.

Per il formato DOCX, invece, il dataset è composto da circa 5\,000 file reali e
2\,000 file randomici.


\subsection{Analisi per Formato}

In questa sezione sono riportati i risultati ottenuti nella valutazione delle performance del detector sui file di test.

Per ogni formato si considera:
\begin{itemize}
    \item il numero e la percentuale di falsi positivi (FP)
    \item il numero e la percentuale di falsi negativi (FN)
    \item il contributo di ciascuna metrica al processo decisionale.
\end{itemize}

I risultati che seguono rappresentano l'aggregazione dei risultati ottenuti sui tre fold.

\begin{table}[H]
\centering
\begin{tabular}{lcccc}
\toprule
Metrica & FN & FN (\%) & FP & FP (\%) \\
\midrule
JSD     & 6  & 0.06 & 0 & 0.00 \\
TVD     & 6  & 0.06 & 0 & 0.00 \\
$L_1$   & 5  & 0.05 & 0 & 0.00 \\
Cosine  & 0  & 0.00 & 1 & 0.05 \\
Entropy & 150 & 1.50 & 11 & 0.59 \\
\bottomrule
\end{tabular}
\caption{Risultati del detector per il formato JPG.}
\end{table}


\begin{table}[H]
\centering
\begin{tabular}{lcccc}
\toprule
Metrica & FN & FN (\%) & FP & FP (\%) \\
\midrule
JSD     & 1 & 0.02 & 0 & 0.00 \\
TVD     & 1 & 0.02 & 0 & 0.00 \\
$L_1$   & 1 & 0.02 & 0 & 0.00 \\
Cosine  & 2 & 0.04 & 0 & 0.00 \\
Entropy & 2 & 0.04 & 0 & 0.00 \\
\bottomrule
\end{tabular}
\caption{Risultati del detector per il formato DOCX.}
\end{table}



\begin{table}[H]
\centering
\begin{tabular}{lcccc}
\toprule
Metrica & FN & FN (\%) & FP & FP (\%) \\
\midrule
JSD     & 0 & 0.00 & 0 & 0.00 \\
TVD     & 0 & 0.00 & 0 & 0.00 \\
$L_1$   & 0 & 0.00 & 0 & 0.00 \\
Cosine  & 0 & 0.00 & 0 & 0.00 \\
Entropy & 3 & 0.03 & 0 & 0.00 \\
\bottomrule
\end{tabular}
\caption{Risultati del detector per il formato PDF.}
\end{table}




\begin{table}[H]
\centering
\begin{tabular}{lcccc}
\toprule
Metrica & FN & FN (\%) & FP & FP (\%) \\
\midrule
JSD     & 116 & 1.16 & 0 & 0.00 \\
TVD     & 139 & 1.39 & 0 & 0.00 \\
$L_1$   & 131 & 1.31 & 0 & 0.00 \\
Cosine  & 73  & 0.73 & 1 & 0.05 \\
Entropy & 8   & 0.08 & 0 & 0.00 \\
\bottomrule
\end{tabular}
\caption{Risultati del detector per il formato TXT.}
\end{table}


Dai risultati, si può osservare che il sistema di rilevamento implementato riesce a classificare correttamente
la maggior parte dei file reali e randomici, con percentuali di falsi positivi prossime allo 0\% per tutti i formati considerati.

\bigskip

Anche le percentuali di falsi negativi risultano molto basse, seppure leggermente più elevate,
con valori generalmente prossimi allo 0\% e comunque inferiori al 2\% per tutti i formati analizzati.


Nel dettaglio, si osservano percentuali di falsi negativi leggermente più elevate
per i formati JPG e TXT. Tale comportamento può essere ricondotto alla presenza,
nei dataset utilizzati, di una piccola percentuale di file con caratteristiche
strutturali atipiche o fortemente entropiche, che risultano più simili a
distribuzioni casuali.

Tuttavia, anche in questi casi, le percentuali di errore rimangono estremamente
ridotte e non compromettono in alcun modo le prestazioni complessive del modello.

\bigskip

Per concludere, questi risultati indicano che il modello è efficace nel distinguere tra file reali e generati casualmente,
dimostrando la validità dell'approccio basato sulle N-gram e sulle metriche di distanza.


\subsection{Valutazione Operativa su Dataset Esteso}

Per valutare ulteriormente l'efficacia del sistema di rilevamento, è stata condotta una valutazione operativa
su un dataset esteso di file reali, non utilizzati nelle fasi di addestramento e validazione.

\bigskip

Nel dettaglio, sono stati considerati:
\begin{itemize}
    \item circa 100\,000 file JPG reali;
    \item circa 80\,000 file TXT reali.
\end{itemize}

\bigskip

Il detector è stato eseguito utilizzando i centroidi e le soglie finali ottenute
dall’aggregazione dei risultati sui diversi fold di cross-validation, senza
ulteriori fasi.

\bigskip

Su 109\,233 file JPG reali, il detector ha classificato come reali il 97,5\% dei file,
con una percentuale di falsi negativi pari al 2,5\%.

L'analisi delle metriche mostra come la maggior parte dei falsi negativi sia dovuta alla metrica di entropia,
la quale ha classificato come non reali 2730 file sui 2749 file totali classificati non reali.

Questo comportamento è coerente con quanto osservato nella fase di validazione, dove la metrica di entropia ha mostrato una maggiore sensibilità
ai file con caratteristiche strutturali atipiche.

\bigskip

Per quanto riguarda i file TXT, su un totale di 78\,285 file reali, il sistema ha classificato come reali il 99,32\% dei file,
con una percentuale di falsi negativi pari allo 0,68\%.

In questo caso tutte le metriche hanno contribuito in modo equilibrato alla classificazione,
senza che una singola metrica emergesse come predominante nella generazione dei falsi negativi.

\bigskip

Questi risultati confermano l'efficacia del sistema di rilevamento anche su dataset estesi e non visti durante l'addestramento,
dimostrando la capacità del modello di generalizzare e mantenere alte prestazioni in scenari operativi reali.




\section{Considerazioni sulle Scelte Progettuali e Confronto con Approcci di Machine Learning}

In fase di progettazione del sistema di rilevamento, si è valutata la possibilità
di adottare tecniche di Machine Learning per la classificazione dei file.

Tuttavia, dopo un'attenta valutazione, si è deciso di optare per un approccio basato su metriche di distanza e soglie ottimizzate,
per diverse ragioni.

\bigskip

In primo luogo, le tecniche di Machine Learning richiedono generalmente grandi quantità di dati per l'addestramento,
e nel contesto di questo lavoro, la disponibilità di dataset ampi e bilanciati per ciascun formato di file era limitata.

\bigskip

Inoltre, l'addestramento di modelli generativi o discriminativi può essere computazionalmente costoso,
richiedendo risorse hardware significative e tempi di addestramento lunghi.

\bigskip

Infine, l'approccio basato su metriche di distanza offre una maggiore trasparenza e interpretabilità,
consentendo di comprendere facilmente come viene presa la decisione di classificazione,
cosa che può essere più complessa con modelli di Machine Learning più sofisticati.

\bigskip

Pertanto, l'approccio adottato rappresenta una soluzione efficace e pragmatica,
bilanciando la necessità di accuratezza con la praticità e l'efficienza computazionale.



\section{Lavori Correlati e Confronto con Magika}

In questi ultimi anni sono stati individuati diversi metodi per la verifica del tipo di file basati su analisi statistiche del 
contenuto binario, indipendentemente dall'estensione. In questo contesto, uno dei lavori più rilevanti è Magika,
un tool sviluppato da Google per l'identificazione del tipo di file basata soltanto sull'analisi statistica del contenuto.

\bigskip

L'obiettivo di Magika è quello di riconoscere in modo affidabile il formato di un file anche se questo 
presenta un offuscamento del contenuto, come la modifica dell’estensione o di parti dell’header.
Questo dimostra come l'identità di un file sia fortemente legata all'insieme di pattern statistici presenti nel suo contenuto binario,
piuttosto che ai singoli byte o sequenze di byte specifiche.

\bigskip

Il presente lavoro affronta un problema affine, tuttavia con obiettivi differenti.
Magika infatti si propone di identificare il formato di un file, mentre l'obiettivo qui è quello di valutare la
verosimiglianza di un file rispetto a un formato atteso, con il fine ultimo di proteggere dati sensibili da minacce ransomware.

\bigskip

Dal punto di vista metodologico, le differenze sono evidenti.
Magika utilizza un approccio basato su modelli di Machine Learning, in particolare reti neurali, per apprendere le caratteristiche distintive 
di ciascun formato di file.
Questo richiede un addestramento su dataset di grandi dimensioni e una decisione di classificazione che può risultare meno interpretabile.
Al contrario, il presente lavoro adotta un approccio basato su rappresentazioni esplicite come le N-gram e su
metriche di distanza ben definite, che consentono una maggiore trasparenza nel processo decisionale.

\bigskip

Benché le differenze fra i due approcci siano significative, si possono comunque fare alcune considerazioni comparative.
Strumenti come Magika dimostrano che l'identità di un file può essere catturata efficacemente attraverso l'analisi statistica del contenuto binario,
risulta quindi rilevante studiare in che modo questa identità possa essere manipolata mantenendo inalterato il contenuto originale.   
In questo senso, il presente lavoro può essere considerato come un passo in più verso 
l'analisi controllata dell'identità statistica di un file, fornendo strumenti utili per la rilevazione e per 
la generazione di file in grado di eludere varie minacce.



\chapter{Sistema di Generazione dei File Verosimili}

Nel capitolo precedente è stato descritto nel dettaglio il sistema di rilevamento implementato,
illustrando le scelte progettuali, il flusso di lavoro e i risultati sperimentali ottenuti.

\bigskip

Il detector così definito non rappresenta esclusivamente uno strumento di valutazione,
ma costituisce un riferimento quantitativo per la validazione delle tecniche di
generazione di file verosimili. Attraverso l'uso di rappresentazioni statistiche
e di metriche di distanza, esso consente di stabilire in modo oggettivo se un file
artificiale rispetti le caratteristiche strutturali tipiche dei file reali per un
determinato formato.

\bigskip

In questo contesto, il sistema di rilevamento verrà utilizzato come oracolo per verificare l'efficacia
delle tecniche di generazione di file verosimili, al fine di
evitare sia una replicazione dei file reali, sia la generazione di file completamente randomici.

In questo capitolo verranno illustrate le metodologie adottate per la generazione di file verosimili,
progettate per risultare statisticamente compatibili con i vincoli imposti dal detector.

\section{Idea Generale}

Come descritto nei capitoli precedenti, per generare file verosimili l'idea è che questi
debbano risultare statisticamente compatibili con i file reali, senza ricorrere
alla replica diretta di file esistenti o alla generazione di file completamente casuali.

\bigskip

L'obiettivo non è quello di generare file semanticamente validi o utilizzabili in applicazioni reali,
ma piuttosto di riprodurre caratteristiche osservabili a livello binario.

In particolare, si mira a preservare le proprietà statistiche di ciascun formato, in modo tale
che i file generati non risultino distinguibili dai file reali attraverso l'analisi statistica.

\bigskip

La strategia adottata prevede una prima modellazione statistica dei file reali, al fine 
di catturare le caratteristiche strutturali tipiche di ciascun formato, seguita 
da un processo di generazione che utilizza queste caratteristiche per produrre nuovi file verosimili.

\bigskip

I dettagli di questo processo verranno illustrati progressivamente nelle sezioni seguenti, a partire dalla modellazione statistica
fino alla descrizione dell'implementazione del sistema di generazione.

\section{Modellazione Statistica dei File Reali}

Per descrivere in modo efficace le caratteristiche strutturali dei file reali, è necessario
andare oltre la semplice distribuzione globale dei byte e considerare anche le dipendenze locali tra i byte consecutivi.
Come discusso in precedenza, la rappresentazione a bigrammi (N-gram con \(n=2\)) fornisce una base solida per catturare 
queste dipendenze, evidenziando pattern tipici di ciascun formato.

\bigskip

L'idea centrale è quella di modellare la generazione dei byte come un processo stocastico a tempo discreto,
in cui ogni byte è rappresentato da una variabile aleatoria discreta.
In questo contesto, a ciascuna posizione \(n\) del file è associata una variabile aleatoria \(X_n\),
la quale assume valori nell'insieme discreto finito \(S = \{0, \dots, 255\}\),
corrispondente all’insieme dei possibili valori di byte.

\bigskip

L'utilizzo dei bigrammi può essere interpretato come una stima delle probabilità condizionate di osservare un certo 
valore di byte in funzione del byte precendente. Questo riflette l'ipotesi che la probabilità di osservare 
un byte non sia indipendente, ma dipenda dal contesto immediatamente precedente.

\bigskip

Tali assunzioni portano naturalmente alla modellazione della sequenza di byte tramite una catena di Markov 
a tempo discreto di ordine uno.
Si richiama pertanto la seguente definizione:

\begin{definition}[Catena di Markov a tempo discreto]
Si dice catena di Markov a tempo discreto una successione di variabili aleatorie
$(X_n)_{n \geq 1}$, definite su uno stesso spazio di probabilità, tali che:
\begin{enumerate}
    \item le variabili aleatorie $X_n$ sono discrete e assumono valori in uno stesso
    insieme discreto $S$, detto spazio degli stati;
    \item vale la proprietà di Markov, ovvero per ogni $i_1, \dots, i_{n-1}, i, j \in S$
    risulta:
    \[
    \mathbb{P}(X_{n+1} = j \mid X_1 = i_1, \dots, X_n = i)
    =
    \mathbb{P}(X_{n+1} = j \mid X_n = i).
    \]
\end{enumerate}
\end{definition}

\bigskip

Nel contesto di questo lavoro è dunque possibile rappresentare la generazione di un file come un 
processo markoviano, in cui la probabilità di osservare un byte
dipenda esclusivamente dal byte immediatamente precedente, ipotesi coerente con la modellazione
basata su bigrammi adottata e che costituisce il fondamento teorico per la costruzione delle
matrici di transizione descritte nella sezione successiva.


\section{Selezione delle Porzioni di File Analizzate}

In maniera del tutto analoga a quanto fatto per il sistema di rilevamento, è necessario definire
le porzioni di file da analizzare per la costruzione delle matrici di transizione.
Come visto in precedenza, nei file reali vi sono sezioni altamente entropiche che non aggiungono
valore informativo alla modellazione statistica del formato, ma che possono invece
introdurre rumore e distorsioni.

\bigskip

In questo contesto è dunque opportuno selezionare accuratamente le porzioni di file su cui 
costruire le catene, evitando di includere sezioni che potrebbero compromettere la qualità della
modellazione.

\bigskip

Nel dettaglio, si adottano le stesse strategie di selezione delle porzioni di file
utilizzate per il sistema di rilevamento, come descritto nella Sezione 3.3.

\bigskip

L’adozione dello stesso criterio di estrazione strutturale sia nella fase di
rilevamento sia nella fase di generazione garantisce inoltre coerenza tra il modello
statistico appreso e il sistema di valutazione, permettendo di utilizzare il
detector come oracolo per la validazione dei file generati.



\section{Costruzione delle Matrici di Transizione}

Data l'idea di modellare la generazione dei byte come un processo markoviano,
il passo successivo consiste nella costruzione delle matrici di transizione che
descrivono le probabilità condizionate di transizione tra i vari stati (byte).

Nel dettaglio, una catena di Markov a tempo discreto è caratterizzata dalle probabilità
di transizione tra i suoi stati.

Nel nostro caso, dove la catena risulta essere omogenea e a stati finiti, 
tali probabilità possono essere rappresentate tramite una matrice di transizione \(P\).

\bigskip

Si ricorre alla seguente definizione:

\begin{definition}[Matrice di transizione]
Sia \((X_n)_{n \geq 1}\) una catena di Markov omogenea a stati finiti, con spazio
degli stati \(S\) di cardinalità \(N\). Si definisce matrice di transizione la
matrice \(\Pi \in \mathbb{R}^{N \times N}\) le cui componenti sono date da:
\[
\pi_{ij} = \mathbb{P}(X_{n+1} = j \mid X_n = i),
\quad \forall\, i,j \in S,
\]
dove l’elemento
\(\pi_{ij}\) rappresenta la probabilità di transizione dallo stato \(i\) allo
stato \(j\).
\end{definition}

\bigskip

In questo lavoro, lo spazio degli stati \(S\) coincide con l'insieme dei possibili valori di byte,
\(S = \{0, \dots, 255\}\).
Ne consegue che la matrice di transizione \(P\) è una matrice \(256 \times 256\) in cui
ogni riga rappresenta la distribuzione di probabilità condizionata del byte successivo dato il byte corrente.

\bigskip

Per costruire la matrice di transizione per ciascun formato di file, si procede come segue:

\bigskip   
\begin{itemize}
    \item Dato un insieme di file reali appartenenti a uno stesso formato, si analizzano le
sequenze di byte estratte dai file e, per ogni coppia di byte consecutivi
\((i,j)\) osservata, viene incrementato il contatore associato alla transizione
dallo stato \(i\) allo stato \(j\).

    \item In questo modo, si ottiene una matrice di conteggi \(C\) in cui ogni elemento
\(C_{ij}\) rappresenta il numero di occorrenze della coppia di byte \((i,j)\) osservata
nei file reali.

     \item Ottenuta la matrice di conteggi, la matrice di transizione \(P\) viene calcolata normalizzando ciascuna riga, 
ovvero:
\[
P_{ij} = \frac{C_{ij}}{\sum_{k \in S} C_{ik}}, \quad i,j \in S,
\]
in modo tale che ogni riga di \(P\) rappresenti una distribuzione di probabilità condizionata valida.

     
\end{itemize}


Questo processo viene ripetuto per ciascun formato di file considerato,
ottenendo così una matrice di transizione specifica per ogni formato.


\section{Processo di Generazione Markoviana}

Una volta ottenuta la matrice di transizione \(P\), il passo successivo consiste nell'utilizzare questa matrice
per generare, byte dopo byte, nuovi file verosimili.

Il processo di generazione viene modellato come un campionamento sequenziale dalla catena di markov
definita dalla matrice di transizione \(P\).
In particolare, definita \(L\) la lunghezza del file da generare, il file viene generato producendo
una sequenza di variabili aleatorie \((X_1, X_2, \dots, X_L)\), dove ciascuna variabile \(X_n\) rappresenta il byte 
alla posizione \(n\) del file.

\bigskip

Il primo byte della sequenza \(X_1\) viene generato casualmente, selezionando uno stato iniziale \(i_0\)
nello spazio degli stati \(S\).

A partire da questo stato iniziale, ogni byte successivo \(X_{n+1}\) viene generato campionando dalla
distribuzione di probabilità condizionata associata allo stato corrente \(X_n\), 
rappresentato dalla riga corrispondente della matrice di transizione \(P\).

\bigskip

Formalmente, definito il byte corrente \(X_n = i\), il byte successivo \(X_{n+1}\) viene generato come:
\[
\mathbb{P}(X_{n+1} = j \mid X_n = i) = P_{ij}, \quad j \in S.
\]
In questo modo ogni transizione tra byte avviene con
una probabilità coerente con quelle apprese durante la fase di costruzione della
matrice di transizione

\bigskip

Dal punto di vista del campionamento, la scelta del byte successivo avviene a partire 
dalla distribuzione discreta definita dalla riga \(P_{i\cdot}\) della
matrice di transizione.

A tal fine, si costruisce la funzione di distribuzione cumulativa (CDF) associata a questa riga, definita come:
\[
F_i(j) = \sum_{k = 0}^{j} P_{ik}, \quad j \in S.
\]

Dato un numero casuale \(r\) estratto uniformemente nell'intervallo \([0, 1)\),
il byte successivo è determinato come il più piccolo indice \(j \in S\) tale che:
\[
F_i(j) \geq r.
\]

\bigskip

Per fare un esempio, si consideri una riga della matrice di transizione
associata allo stato corrente \(i\), limitata a un sottoinsieme di stati
\(S = \{0,1,2,3\}\), con probabilità:
\[
P_{i\cdot} = (0.1,\; 0.3,\; 0.4,\; 0.2).
\]
La corrispondente funzione di distribuzione cumulativa risulta quindi:
\[
F_i = (0.1,\; 0.4,\; 0.8,\; 1.0).
\]
Supponendo di estrarre un numero casuale \(r = 0.35\), il byte successivo è determinato come 
il più piccolo indice \(j\) tale che \(F_i(j) \geq 0.35\).
In questo caso, si ha \(F_i(1) = 0.4 \geq 0.35\), pertanto il byte successivo sarà \(X_{n+1} = 1\).



\section{Ottimizzazione del Campionamento}

Come descritto nella sezione precedente, il campionamento diretto da una distribuzione discreta
può essere effettuato costruendo la funzione di distribuzione cumulativa e selezionando il primo
stato che soddisfa la condizione \(F_i(j) \geq r\).

Tale approccio, sebbene corretto, richiede una complessità lineare \(O(N)\) rispetto al numero di stati \(N\).

\bigskip

In questo contesto, lo spazio degli stati \(S\) ha cardinalità pari al numero di possibili valori di byte,
ovvero \(S = 256\). Di conseguenza, il campionamento diretto non risulta computazionalmente costoso sul singolo file,
tuttavia considerando l'elevato numero di file da generare e la lunghezza di ciascun file,
questa operazione può diventare un collo di bottiglia in termini di prestazioni.
Pertanto l'ottimizzazione di questa fase può portare a significativi risparmi in termini di tempo di esecuzione complessivo.

\bigskip

Per risolvere questa problematica, si è scelto di ricorrere all’Alias Method, una
tecnica che consente di effettuare il campionamento in tempo costante \(O(1)\),
una volta eseguita una fase di preprocessamento.

\bigskip

Il metodo si basa sulla trasformazione della distribuzione discreta originale in
una rappresentazione equivalente che utilizza due strutture dati ausiliarie: un
array di probabilità \texttt{prob} e un array di alias \texttt{alias}.
Nel dettaglio, dato un insieme di \(N\) stati con probabilità associate \(p_1, p_2, \dots, p_N\),
con \(\sum_{i=1}^{N} p_i = 1\), l’Alias Method costruisce le tabelle \texttt{prob} e \texttt{alias}
che permettono di campionare dalla distribuzione in tempo costante.

\bigskip

In fase di preprocessamento, le probabilità sono moltiplicate per \(N\), ottenendo
una distribuzione scalata \(q_i = N \cdot p_i\) \cite{darago_alias_2021}.
Successivamente, gli stati vengono suddivisi in due insiemi:
\begin{itemize}
  \item stati con \(q_i < 1\), detti \textit{small};
  \item stati con \(q_i \ge 1\), detti \textit{large}.
\end{itemize}

\bigskip


Si procede poi associando iterativamente uno stato \(s \in \textit{small}\) a uno stato \(l \in \textit{large}\),
assegnando \(\texttt{prob}[s] = q_s\) e \(\texttt{alias}[s] = l\), e aggiornando
\(q_l \leftarrow q_l - (1 - q_s) = q_l + q_s - 1\).

Il processo viene ripetuto fino a esaurimento degli insiemi; gli stati rimanenti vengono completati
impostando \(\texttt{prob}[i]=1\) e \(\texttt{alias}[i]=i\), garantendo la correttezza della rappresentazione.

\bigskip

Una volta costruite le tabelle, il campionamento da una distribuzione discreta può essere effettuato
in tempo costante: si seleziona uniformemente un indice \(k \in \{1, \dots, N\}\) e si estrae un numero casuale
\(r \in [0,1)\). Se \(r < \texttt{prob}[k]\), lo stato selezionato è \(k\); altrimenti viene restituito lo stato
\(\texttt{alias}[k]\) \cite{walker1977alias}.

\bigskip

Questo procedimento garantisce che la distribuzione degli stati campionati coincida con quella originale,
pur richiedendo un numero costante di operazioni per ogni estrazione.




\section{Implementazione del Sistema di Generazione}

In questa sezione vengono descritti gli aspetti implementativi del sistema di generazione dei file.
L’obiettivo è tradurre il modello teorico
presentato nelle sezioni precedenti in una pipeline concreta ed efficiente,
in grado di produrre grandi quantità di file sintetici.

\bigskip

L'implementazione sviluppata segue la struttura generale delineata nelle sezioni precedenti:
la generazione di file verosimili avviene attraverso un processo markoviano basato su matrici di transizione
costruite a partire da file reali, con l'ottimizzazione del campionamento tramite l'Alias Method.


\subsection{Struttura Generale}

Il sistema di generazione è organizzato come una pipeline modulare, suddivisa in due fasi principali:
una fase di preprocessamento \emph{offline}, la quale si occupa della costruzione delle matrici di transizione, e una
fase di generazione \emph{online}, responsabile della produzione dei file.

\bigskip

La separazione tra le due fasi consente di ottimizzare il processo,
permettendo di isolare le operazioni computazionalmente costose nella fase di preprocessamento,
al fine di migliorare le prestazioni del sistema.

\bigskip

Come anticipato, la fase offline ha come obiettivo l'apprendimento delle distribuzioni di transizione tra byte
a partire da file reali. Per ciascun formato di file considerato,
vengono analizzati un insieme di file reali dai quali sono estratte le porzioni di file discusse in precedenza.

\bigskip

A partire da queste porzioni, vengono costruite le matrici di transizione
secondo la metodologia descritta nella Sezione 5.4, seguita dalla fase di preprocessamento
per l'Alias Method, come descritto nella Sezione 5.6.

Il risultato di questa fase consiste in strutture dati ausiliarie che
permettono di effettuare il campionamento delle transizioni in tempo costante
durante la generazione runtime.

\bigskip

Per quanto riguarda la fase online, il sistema utilizza esclusivamente i modelli discussi in precedenza.

Dato un numero prefissato di file da generare, si genera per ciascun file una sequenza di byte
di lunghezza variabile, campionando iterativamente dalla catena di Markov mediante le tabelle alias.

Inoltre, fissato N il numero di file da generare, si generano $\frac{N}{4}$  file per ciascun formato considerato,
ottenendo così una distribuzione bilanciata tra i vari formati.

\bigskip

La generazione è indipendente dal dataset originale e non richiede ulteriori
operazioni di analisi sui file reali, rendendo il processo efficiente e facilmente
scalabile. Le sequenze di byte generate vengono infine incapsulate nel formato
di output desiderato, secondo le modalità specifiche di ciascun formato.

\bigskip

L’intera pipeline viene applicata in maniera
analoga a tutti i formati considerati (PDF, JPG, DOCX e TXT). Le differenze tra i
formati riguardano esclusivamente le modalità di estrazione delle porzioni
strutturali e la fase di ricostruzione finale del file, mentre il nucleo del
processo di generazione rimane invariato.


\subsection{Costruzione delle Matrici di Markov}

Come descritto precedentemente, la costruzione delle matrici di transizione avviene
nella fase di preprocessamento offline.

\bigskip

Per prima cosa, è necessario costruire le matrici di conteggi a partire dai file reali.

In tal senso si presenta l'algoritmo seguente:

\bigskip

\begin{algorithm}[H]
\caption{countBigrams}

\KwIn{matrice $C$, sequenza di byte $B$}


\For{$n \gets 1$ \KwTo $|B| - 1$}{
    $i \gets B[n]$\;
    $j \gets B[n+1]$\;
    $C_{ij} \gets C_{ij} + 1$\;
}

\Return $C$\;
\end{algorithm}

\bigskip

L'algoritmo \texttt{countBigrams} prende in input una sequenza di byte \(B\), e aggiorna la matrice \(C\) incrementando i contatori
corrispondenti alle coppie di byte osservate nella sequenza.

\bigskip

Una volta costruita la matrice di conteggi globale \(C\), si procede alla
costruzione della matrice di transizione \(P\) normalizzando ciascuna riga.

\bigskip    

\begin{algorithm}[H]
\caption{normalizeRows}

\KwIn{Matrice di conteggi $C \in \mathbb{N}^{256 \times 256}$}

\For{$i \gets 0$ \KwTo $255$}{
    $s \gets \sum_{j=0}^{255} C_{ij}$\;
    \eIf{$s > 0$}{
        \For{$j \gets 0$ \KwTo $255$}{
            $P_{ij} \gets \dfrac{C_{ij}}{s}$\;
        }
    }{
        \For{$j \gets 0$ \KwTo $255$}{
            $P_{ij} \gets \dfrac{1}{256}$\;
        }
    }
}

\Return $P$\;
\end{algorithm}

\bigskip

La matrice di transizione così ottenuta rappresenta una stima delle
probabilità di transizione della catena di Markov associata al formato
considerato, ed è successivamente utilizzata per la fase di
preprocessamento (costruzione delle tabelle alias).

\bigskip

Implementati gli algoritmi sopra descritti, si procede con la costruzione delle matrici di transizione
per ciascun formato di file considerato, utilizzando i dataset di file reali
descritti nel Capitolo~4.

\bigskip

\begin{algorithm}[H]
\caption{BuildTransitionMatrix}

\KwIn{Dataset di file reali $\mathcal{D}$}

Inizializza $C \gets \mathbf{0}_{256 \times 256}$\;

\ForEach{file $f \in \mathcal{D}$}{
    $B \gets$ estrai\_byte\_strutturali($f$)\;
    \If{$|B| \geq 2$}{
        $C \gets$ countBigrams($C, B$)\;
    }
}

$P \gets$ normalizeRows($C$)\;

\Return $P$\;
\end{algorithm}



\subsection{Preprocessamento tramite Alias Method}

Come descritto nella sezione precedente, l’Alias Method consente di effettuare
il campionamento da una distribuzione discreta in tempo costante, a fronte di
una fase di preprocessamento.

\bigskip

Tale preprocessamento viene applicato alle
matrici di transizione markoviane costruite nella fase precedente. Poiché ogni
riga della matrice di transizione \(P\) rappresenta una distribuzione di
probabilità discreta sullo spazio degli stati \(S\), è possibile costruire una
tabella alias indipendente per ciascuna riga.

\bigskip

In dettaglio, il preprocessamento produce due matrici ausiliarie:
\[
\text{Prob} \in [0,1]^{256 \times 256}, \qquad
\text{Alias} \in \{0,\dots,255\}^{256 \times 256},
\]
tali che la riga \(i\) di ciascuna matrice consente di campionare il byte
successivo dato lo stato corrente \(i\).

\bigskip

\begin{algorithm}[H]
\caption{PreprocessingAliasMethod}

\KwIn{Matrice di transizione $P \in [0,1]^{256 \times 256}$}

\For{$i \gets 0$ \KwTo $255$}{
    $(\text{Prob}_{i\cdot}, \text{Alias}_{i\cdot})
    \gets \text{AliasMethod}(P_{i\cdot})$\;
}

\Return $\text{Prob}, \text{Alias}$\;
\end{algorithm}

\subsection{Gestione della Lunghezza dei File}

Un ulteriore aspetto da considerare, prima di concentrare l'attenzione sulla generazione runtime,
riguarda la gestione della lunghezza dei file generati.

Questa infatti non è fissa ma si è preferito
adottare una strategia che consenta di variare la lunghezza dei file generati, al fine di
ottenere un insieme di file sintetici eterogeneo e, al contempo, mantenere
contenuto il costo computazionale della fase di generazione.

\bigskip

Poiché il costo della generazione markoviana cresce linearmente con la lunghezza
del file, una distribuzione uniforme delle lunghezze tenderebbe a sovrarappresentare
file di grandi dimensioni, aumentando significativamente il tempo di esecuzione
complessivo senza apportare un beneficio in termini di informazione
strutturale.

\bigskip

Per questo motivo, si è scelto di adottare una distribuzione log-uniforme all'interno 
di un intervallo \([L_{\min}, L_{\max}]\), in modo da favorire la generazione di file più corti,
pur consentendo la presenza di file di dimensioni maggiori.

\bigskip

Formalmente, la lunghezza \(L\) viene generata come:
\[
L = \left\lfloor \exp(U) \right\rfloor,
\quad \text{con } U \sim \mathcal{U}(\log L_{\min}, \log L_{\max}),
\]
dove \(\mathcal{U}(a,b)\) denota la distribuzione uniforme continua
nell'intervallo \([a,b]\).


\subsection{Scrittura dei File su Disco}

Un ultimo aspetto implementativo prima di definire il generatore riguarda la scrittura dei file generati su disco.

\bigskip

In dettaglio, per i file di tipo TXT, la sequenza di byte generata viene scritta direttamente
su disco, in quanto il formato non impone vincoli strutturali ulteriori.

Analogamente, per i file PDF e JPG, i byte generati vengono scritti direttamente
in un file binario. In questi casi, la generazione mira a preservare 
la validità sintattica intervenendo solo su porzioni selezionate del formato.

\bigskip

Per quanto riguarda i file DOCX, la scrittura su disco richiede un passaggio
aggiuntivo. Infatti, un file DOCX è in realtà un archivio ZIP contenente diversi
file XML e risorse correlate. 
Pertanto, è necessario l'utilizzo di un template DOCX predefinito, nel quale si mantengono i file 
ausiliari e si sostituiscono esclusivamente i file XML contenenti il testo con la sequenza di byte generata.

\bigskip

Nel dettaglio, in fase di generazione le sequenze di byte prodotte dal modello markoviano
vengono inserite esclusivamente all'interno dei file XML principali
(\texttt{document.xml} e \texttt{styles.xml}), sostituendo il contenuto
corrispondente nel template. Il file DOCX finale viene quindi ricostruito
scrivendo l'archivio ZIP risultante su disco.

\bigskip

Questo approccio preserva la struttura del template e vincola la generazione alle sole porzioni di contenuto selezionate,
riducendo il rischio di corruzione del file.



\subsection{Generazione Runtime}

Una volta completata la fase di preprocessamento, è possibile procedere con la generazione runtime
dei file, la quale si basa sul campionamento iterativo dalla catena di Markov
utilizzando le tabelle alias costruite in precedenza.

\bigskip

Il campionamento di un singolo stato a partire dalle tabelle alias e la
generazione iterativa della sequenza di byte vengono formalizzati nei seguenti
algoritmi, utilizzati durante la fase di generazione runtime.

\bigskip

\begin{algorithm}[H]
\caption{sampleAlias}

\KwIn{Vettori $\text{prob}$, $\text{alias}$}

$i \sim \{0,\dots,255\}$\;
$r \sim \mathcal{U}(0,1)$\;

\eIf{$r < \text{prob}[i]$}{
    \Return $i$\;
}{
    \Return $\text{alias}[i]$\;
}
\end{algorithm}

\bigskip

\begin{algorithm}[H]
\caption{MarkovGenerate}

\KwIn{Matrici $\text{Prob}$, $\text{Alias}$, lunghezza $L$, stato iniziale $y_1 \in S$}
$B[1] \gets y_1$\;

\For{$n \gets 2$ \KwTo $L$}{
    $y_n \gets$ sampleAlias($\text{Prob}_{y_{n-1}\cdot}$, $\text{Alias}_{y_{n-1}\cdot}$)\;
    $B[n] \gets y_n$\;
}

\Return $B$\;
\end{algorithm}

\bigskip

Sulla base delle componenti descritte in precedenza, il processo
completo di generazione runtime può essere riassunto
nell’algoritmo seguente.

\bigskip

\begin{algorithm}[H]
\caption{GenerateFiles}
\KwIn{
$N$,
$\text{Prob}_f$, $\text{Alias}_f$,
templateDOCX,
$[L_f^{\min}, L_f^{\max}]$
}
\ForEach{formato $f$}{
    \For{$k \gets 1$ \KwTo $N_f$}{
        $L \sim \text{LogUniform}(L_f^{\min}, L_f^{\max})$\;
        $y_1 \sim \mathcal{U}\big(\{0,\dots,255\}\big)$\;

        $B \gets$ MarkovGenerate($\text{Prob}_f$, $\text{Alias}_f$, $L$, $y_1$)\;

        \eIf{$f = \texttt{DOCX}$}{
            writeFile($B$, $f$, templateDOCX)\;
        }{
            writeFile($B$, $f$)\;
        }
    }
}
\end{algorithm}


\section{Analisi del Costo Computazionale}

In questa sezione viene analizzato il costo computazionale del sistema di
generazione dei file, distinguendo tra la fase di preprocessamento
offline e la fase di generazione runtime.

\bigskip

Poiché la fase offline viene eseguita una sola volta e non è soggetta a vincoli
di tempo, l’analisi si concentra sul costo di memoria
associato ai modelli generati.


\subsection{Fase di Preprocessamento Offline}

Per ciascun formato di file, si memorizzano due matrici ausiliarie per l’Alias Method, $\text{Prob}$ e $\text{Alias}$,
entrambe di dimensione $256 \times 256$.

\bigskip

Il costo di memoria è pertanto dominato da un numero costante di matrici, ognuna delle quali richiede
$256 \times 256$ elementi, indipendentemente dal numero di file reali utilizzati per la costruzione del modello.

\bigskip

Dal punto di vista asintotico, il costo di memoria per ciascun formato risulta
pari a $O(|S|^2)$, con $|S| = 256$.

\subsection{Fase di Generazione Runtime}

In questa fase, l'analisi si concentra sul costo temporale per la generazione dei file.

\bigskip

La generazione si basa sugli algoritmi \texttt{sampleAlias} e \texttt{MarkovGenerate} descritti in precedenza.

\bigskip

L'algoritmo \texttt{sampleAlias} esegue un numero costante di operazioni, 
pertanto il costo temporale di questo algoritmo è \(O(1)\).

\bigskip

L'algoritmo \texttt{MarkovGenerate}, invece, genera una sequenza di byte di lunghezza \(L\),
invocando l'algoritmo \texttt{sampleAlias} una volta per ciascun byte successivo.

Di conseguenza, il costo temporale complessivo per la generazione di una singola sequenza
di byte di lunghezza \(L\) risulta essere:
\[
O(L).
\]

\bigskip

Per ciascun file generato, il costo totale dipende da:
\begin{itemize}
    \item generazione della sequenza di byte tramite la catena di Markov, con costo \(O(L)\);
    \item scrittura della sequenza di byte su disco.
\end{itemize}

Per i formati TXT, PDF e JPG, la scrittura su disco consiste nella scrittura
diretta di \(L\) byte, e ha pertanto costo \(O(L)\).
Il costo totale per un singolo file di questi formati risulta quindi:
\[
O(L).
\]

\bigskip

Per i file DOCX, la scrittura su disco richiede la ricostruzione di un archivio ZIP
a partire da un template predefinito.
Tale operazione comporta:
(i) la scrittura delle porzioni generate (file XML principali), con costo \(O(L)\),
e (ii) la riscrittura dei file statici del template, con costo \(O(S_{\text{tpl}})\),
dove \(S_{\text{tpl}}\) è la dimensione del template DOCX, costante per il formato.
Di conseguenza, il costo complessivo per la generazione di un singolo file DOCX
risulta pari a:
\[
O(L + S_{\text{tpl}}),
\]
che asintoticamente è lineare in \(L\).

\bigskip

Poiché, per ciascun file, sia la fase di generazione sia la fase
di scrittura su disco richiedono un tempo lineare nella lunghezza del file,
il costo complessivo per singolo file è dato dalla somma:
\[
O(L) + O(L) = O(L).
\]


Date le considerazioni sopra esposte, si procede con l'analisi del costo complessivo.

Sia \(N = \sum_f N_f\) il numero totale di file generati e sia
\[
L = \frac{1}{N} \sum_f \sum_{k=1}^{N_f} L_k
\]
la lunghezza media dei file.
Allora il costo totale della fase di generazione runtime può essere riscritto come:
\[
O\!\left( \sum_f \sum_{k=1}^{N_f} L_k \right)
=
O(N \cdot L).
\]





\section{Valutazione dei File Generati}

Una volta descritto nel dettaglio il sistema di generazione dei file verosimili,
è necessario valutare la qualità dei file prodotti.

\bigskip

L'obiettivo di questa fase è verificare in che misura i file generati
con il modello markoviano risultino statisticamente distanti dai file completamente casuali, e, al contempo,
sufficientemente simili ai file reali.

\subsection{Distanza tra File Generati e File Random}

Al fine di valutare se il metodo di generazione discusso produca file effettivamente 
distinguibili da sequenze casuali di byte, è stata condotta un'analisi analoga a quella descritta
nella Sezione~4.1, sostituendo i file reali con i file generati dal modello.

\bigskip

In particolare, l’obiettivo di questa analisi è verificare che i file generati
non presentino le caratteristiche tipiche dei file completamente randomici,
ma mantengano una struttura statistica coerente con il formato di riferimento.
Per garantire coerenza metodologica, vengono adottate le stesse rappresentazioni
(N-gram con $n=2$) e le medesime metriche di distanza e divergenza introdotte
nel Capitolo~3 e già utilizzate per il confronto tra file reali e file random.

\bigskip

Per ciascun formato considerato il procedimento adottato per l'analisi è analogo a quello descritto
nella Sezione~4.1. 

\bigskip    

Nelle figure sono riportati i risultati ottenuti per ciascun formato.

\bigskip

\bigskip

{\centering
\includegraphics[width=0.8\textwidth]{jsd_distance_generated_vs_random.pdf}
\captionof{figure}{Distanza media tra file generati e file random calcolata mediante la JSD.}
\label{fig:jsd_generated_random}
\par}

\bigskip

{\centering
\includegraphics[width=0.8\textwidth]{l1_distance_generated_vs_random.pdf}
\captionof{figure}{Distanza media tra file generati e file random calcolata mediante la Manhattan Distance.}
\label{fig:l1_generated_random}
\par}

\bigskip

{\centering
\includegraphics[width=0.8\textwidth]{tvd_distance_generated_vs_random.pdf}
\captionof{figure}{Distanza media tra file generati e file random calcolata mediante la TVD.}
\label{fig:tvd_generated_random}
\par}

\bigskip

{\centering
\includegraphics[width=0.8\textwidth]{cosine_distance_generated_vs_random.pdf}
\captionof{figure}{Distanza media tra file generati e file random calcolata mediante la Cosine Similarity.}
\label{fig:cosine_generated_random}
\par}

\bigskip

{\centering
\includegraphics[width=0.8\textwidth]{entropy_real_vs_random.pdf}
\captionof{figure}{Entropia media tra file generati e file random.}
\label{fig:entropy_generated_random}
\par}


\bigskip

Coerentemente con quanto fatto nella Sezione 4.1, anche in questo caso i valori degli istogrammi
risultano concentrati in prossimità del massimo e le differenze tra i formati sono poco evidenti a livello visivo.
Per facilitare il confronto diretto con i risultati della Sezione 4.1, si riporta quindi una tabella
con i valori medi delle metriche per ciascun formato


\begin{table}[H]
\centering
\caption{Statistiche medie per formato (generati-random)}
\label{tab:mean-generated-random}
\renewcommand{\arraystretch}{1.25}
\setlength{\tabcolsep}{8pt}
\begin{tabular}{lcccccc}
\toprule
\textbf{Formato} & \textbf{JSD} & \textbf{TVD} & \textbf{L1} & \textbf{Cosine\_sim} & \textbf{Entropy\_random} & \textbf{Entropy\_real} \\
\midrule
docx & 0.9831 & 0.9900 & 1.9800 & 0.0124 & 11.9378 & 8.2319 \\
jpg  & 0.9812 & 0.9945 & 1.9891 & 0.0222 & 14.4568 & 7.5339 \\
pdf  & 0.9739 & 0.9923 & 1.9846 & 0.0240 & 15.2412 & 7.7467 \\
txt  & 0.9631 & 0.9894 & 1.9788 & 0.0330 & 15.5626 & 8.5213 \\
\bottomrule
\end{tabular}
\end{table}




\bigskip

Dai risultati ottenuti emerge come i file generati risultino statisticamente distinti
dai file completamente randomici per tutti i formati considerati e per tutte le metriche adottate.

\bigskip

Si osserva come le distanze medie tra file generati e file random 
siano simili a quelle osservate tra file reali e file random,
come evidenziato dal confronto con i risultati della Sezione~4.1, 
suggerendo che i file generati mantengano una struttura statistica coerente con il formato di riferimento.

\bigskip

In dettaglio, le metriche di distanza e divergenza (JSD, TVD e $L_1$) assumono
valori medi molto elevati per ciascun formato, con valori di JSD compresi
approssimativamente tra 0.96 e 0.98 e valori di Manhattan distance prossimi al massimo
teorico.

\bigskip

Analogamente, la Cosine Similarity mostra valori attorno allo zero per tutti i formati.

\bigskip

Anche l'analisi dell'entropia media conferma questa tendenza, con valori
significativamente inferiori a quelli dei file randomici, e simili a quelli dei file reali.

\bigskip

In conclusione, questi risultati dimostrano che il metodo di generazione proposto è in
grado di produrre file statisticamente distinguibili da sequenze casuali di byte.
Questa evidenza costituisce un primo passo fondamentale verso la
validazione della verosimiglianza statistica dei file generati, che verrà ulteriormente
analizzata nella sezione successiva attraverso la valutazione tramite il detector.

\subsection{Valutazione tramite il Detector}

Dopo aver verificato che i file generati risultino distinti dai file completamente randomici,
il prossimo passo consiste nel valutare se tali file possano essere considerati verosimili
dal detector sviluppato nel Capitolo~4.

\bigskip

In questa fase, il detector viene utilizzato come un \emph{oracolo di verosimiglianza}:
un file viene considerato verosimile se il detector lo classifica come reale,
secondo i criteri definiti in precedenza.

In questo senso, la valutazione tramite il detector consente di stabilire in modo oggettivo se i
file prodotti dal modello di generazione risultino compatibili con la struttura
statistica appresa dai file reali.

\bigskip

La valutazione è stata condotta utilizzando esclusivamente i centroidi e le
soglie finali ottenute dall’aggregazione dei risultati di cross-validation,
senza introdurre ulteriori fasi di addestramento o di ottimizzazione. 

\bigskip

In dettaglio, sono stati generati 100\,000 file artificiali tramite il modello markoviano,
suddivisi equamente tra i quattro formati considerati.

Di questi, il detector ha classificato come reali il 96.80\% (96\,797) dei file totali, 
mentre il 3,20\% (3203) è stato ritenuto non reali.

\bigskip

L’analisi non è stata condotta su una singola istanza del dataset generato.
Al fine di valutare la stabilità dei risultati, l’esperimento è stato ripetuto
più volte, generando in modo indipendente dataset di dimensioni differenti.
In tutti i casi, le percentuali di file classificati come reali e la
distribuzione delle metriche responsabili delle esclusioni hanno mostrato un
comportamento sostanzialmente invariato.


\bigskip

L'analisi delle metriche responsabili della classificazione evidenzia come la
maggior parte delle esclusioni è dovuta alla metrica di entropia, che da sola
contribuisce alla classificazione come non reali di 3004 file.

Le altre metriche contribuiscono in misura minore:
la metrica JSD esclude 189 file, la TVD 117 file, la Manhattan Distance 121 file
mentre la Cosine Similarity 42 file.

\bigskip

Questo comportamento è coerente con quanto osservato nelle fasi di validazione
del detector: la metrica di entropia risulta la più restrittiva e sensibile a
variazioni nella distribuzione globale dei byte, soprattutto per quanto riguarda il formato JPG,
mentre le metriche di distanza mostrano una maggiore tolleranza nei confronti delle variazioni locali
introdotte dal processo di generazione.

\bigskip

Nel complesso, i risultati ottenuti indicano come il modello di generazione proposto sia in grado
di produrre file che, nella maggior parte dei casi, soddisfino i criteri di verosimiglianza
definiti dal detector.

La bassa percentuale di file considerati non reali suggerisce che il processo markoviano riesca a
riprodurre in modo efficace le caratteristiche statistiche fondamentali dei file
reali, evitando al contempo una replica diretta dei dati di partenza.


\subsection{Valutazione tramite Google Magika}

Magika, come descritto nella Sezione 4.9, è uno strumento basato sull'analisi statistica del contenuto binario
dei file, sviluppato per l'identificazione del formato di un file indipendentemente dalla sua estensione. 
Secondo Google, Magika presenta un'accuratezza del 100\% nel riconoscere i formati di file.

Nel contesto di questo lavoro, Magika è stato utilizzato come strumento aggiuntivo al sistema di rilevamento
per valutare la verosimiglianza dei file generati: se Magika è in grado di identificare correttamente l'estensione
di un file generato, allora il file rispetta le caratteristiche statistiche del formato di riferimento.

\bigskip

Per eseguire questa valutazione, sono stati condotti diversi esperimenti:
è stato valutato come Magika si comporti sui file reali (quelli utilizzati per i dataset), sui file 
completamente randomici generati da Ranflood, e sui file generati dal sistema presentato in questo lavoro.

\bigskip

I risultati ottenuti sono i seguenti:

\bigskip

Per quanto riguarda i file reali, questi sono stati classificati correttamente come appartenenti al formato previsto 
nella quasi totalità dei casi.

\bigskip

I file generati tramite Ranflood (circa 1800 per formato), sono stati invece classificati tutti come *Format file UNKNOWN*: infatti nella totalità dei casi
Magika non è riuscito a classificare correttamente i formati dei file processati.

\bigskip

Nel caso dei file generati dal sistema, l'esperimento è stato condotto utilizzando 100\,000 file, con 25\,000 file per ciascun formato
considerato in questo lavoro.

\bigskip


I risultati ottenuti sono riportati nel seguente grafico:

{\centering
\includegraphics[width=0.8\textwidth]{grafico_torta.png}
\captionof{figure}{Classificazione dei file generati tramite Magika.}
\label{fig:magika_torta_generati}
\par}


\bigskip


Dai risultati emerge che, per i formati DOCX e TXT, i file generati riescono a
"ingannare" Magika, venendo correttamente classificati come appartenenti
al formato corretto nella totalità dei casi.

\bigskip

Per quanto riguarda i file TXT, si è scelto di includere nella categoria testuale
anche i file classificati come CSV. Durante l’esperimento, infatti, si è osservato
come numerosi file TXT vengano classificati da Magika come file CSV, non solo nel
caso dei file generati, ma anche nella valutazione dei file reali.

\bigskip

Questo comportamento è attribuibile alle caratteristiche strutturali condivise
tra i due formati: molti file di testo presentano una distribuzione ricorrente di
virgole, punti e separatori, elementi tipici anche dei file CSV. Di conseguenza,
dal punto di vista statistico, le firme binarie dei file TXT e CSV risultano spesso
difficili da distinguere in modo netto.

\bigskip

Per questo motivo, sia nell’analisi dei file reali sia in quella dei file generati,
si è scelto di considerare congiuntamente le classificazioni TXT e CSV come
appartenenti alla stessa categoria testuale.

\bigskip

Inoltre, con riferimento ai file TXT/CSV, dai risultati emerge come il numero di
file classificati in questa categoria superi i 25\,000 elementi. Questo suggerisce
che anche file appartenenti ad altri formati vengano talvolta ricondotti a una
struttura testuale dal classificatore.

\bigskip

Tale comportamento è imputabile alla natura intrinsecamente più flessibile e meno
vincolata dei formati testuali, che possono essere approssimati da sequenze di
byte con distribuzioni statistiche relativamente semplici. In presenza di contenuti
privi di header rigidi o di strutture binarie fortemente caratterizzanti, Magika
tende quindi a privilegiare una classificazione di tipo testuale.

\bigskip

Per quanto riguarda i formati PDF e JPG, su 25\,000 file generati, Magika ha considerato corretti
circa 8\,500 file per ciascun formato. 

Tale risultato è coerente con la maggiore complessità strutturale dei formati,
che presentano vincoli sintattici più rigidi rispetto ai formati testuali e documentali.
In particolare, entrambi i formati sono caratterizzati dalla presenza di header ben definiti,
marker strutturali obbligatori e sezioni binarie ad alta entropia, elementi che non vengono esplicitamente
modellati dal processo di generazione adottato in questo lavoro, basato esclusivamente su proprietà statistiche locali dei byte.

\bigskip

Di conseguenza, sebbene i file generati riescano in una frazione significativa dei casi a riprodurre firme statistiche compatibili
con i formati di riferimento, l’assenza di una struttura sintattica pienamente conforme limita la capacità del
classificatore di riconoscerli in modo affidabile. Il fatto che circa un terzo dei file venga comunque classificato
correttamente indica tuttavia che il modello di generazione cattura alcune regolarità statistiche rilevanti,
pur non garantendo una riproduzione completa della semantica di formato.

\bigskip

I risultati sopra riportati si riferiscono al modello di Magika con la modalità \emph{ad alta confidenza} (high-confidence).

In questa modalità, Magika fornisce previsioni altamente precise ma con un richiamo (recall) inferiore, 
restituendo una classificazione solo quando il punteggio di confidenza supera soglie conservative 
(altrimenti può fornire un’etichetta generica).

\bigskip

Tuttavia, Magika offre una modalità best-guess, per la quale è possibile ottenere un richiamo maggiore,
accettando anche previsioni con un punteggio di confidenza inferiore. In entrambi i casi, Magika associa
a ogni previsione un punteggio di confidenza compreso tra 0 e 1 \cite{google_magika_prediction_modes}

\bigskip

Nel caso in cui si considerassero anche i risultati ottenuti con la modalità best-guess, i file generati
mostrerebbero percentuali di riconoscimento corretto superiori all'80\% per i formati JPG e PDF, migliorando 
significativamente rispetto alla modalità high-confidence.

Per quanto riguarda i file completamente random invece, anche utilizzando questa modalità i risultati 
rimangono invariati, in quanto neanche con un livello di confidenza più basso Magika riesce a 
valutare correttamente i file.

\bigskip

Per concludere, il confronto con Magika evidenzia come il modello di generazione proposto sia in grado di riprodurre 
file che, pur non essendo identici ai file reali, presentano caratteristiche statistiche sufficientemente simili 
da essere riconosciuti correttamente in molti casi, soprattutto per i formati DOCX e TXT. Questo supporta
l'efficacia del modello di generazione nel creare file verosimili.


\subsection{Valutazione delle Prestazioni: Magika vs Detector}

Un aspetto interessante da considerare riguarda il confronto tra le prestazioni
del detector sviluppato in questo lavoro e quelle di Magika.

Questo confronto è particolarmente rilevante per determinare l’operabilità
del sistema, soprattutto in scenari reali con grandi quantità di file.

\bigskip

Il confronto è stato effettuato in termini di tempo di elaborazione per file.
Per Magika, è stata misurata la durata media del processo di classificazione di un
file, mentre per il detector è stato misurato il tempo di calcolo delle metriche.

\bigskip

\begin{table}[htbp]
    \centering
    \begin{tabular}{|c|c|c|}
    \hline
    \textbf{Strumento} & \textbf{Tempo medio per file} & \textbf{Numero di file processati} \\
    \hline
    Magika & 0.004490 secondi & 100\,000 \\
    Detector & 0.001290 secondi & 100\,000 \\
    \hline
    \end{tabular}
    \caption{Confronto della velocità tra Magika e il detector.}
    \label{tab:performance_comparison}
\end{table}

\bigskip

Dai risultati ottenuti, emerge una differenza significativa in termini di
prestazioni tra Magika e il detector sviluppato in questo lavoro.
In particolare, Magika richiede in media 0.004490 secondi per la classificazione
di un singolo file, mentre il detector presenta un tempo medio pari a 0.001290
secondi per file.

\bigskip

Considerando un dataset di 100\,000 file, questo si traduce in un tempo di
elaborazione complessivo pari a circa 449 secondi (circa 7 minuti)
per Magika, contro circa 129 secondi (circa 2 minuti) per il detector.

\bigskip

Questa differenza è coerente con la natura dei due approcci.
Magika si basa su modelli di Machine Learning profondi, che richiedono un processo
di inferenza più complesso per ciascun file, mentre il detector proposto utilizza
rappresentazioni esplicite basate su N-gram e metriche di distanza computazionalmente
più leggere.

\bigskip

È importante sottolineare che l’obiettivo di questo confronto non è quello di
valutare quale strumento sia complessivamente migliore, ma di evidenziare come
il detector sviluppato in questo lavoro risulti particolarmente adatto a scenari
operativi che richiedono l’elaborazione di grandi quantità di file in tempi
contenuti.

\bigskip

Nel complesso, il confronto delle prestazioni mostra che il detector proposto
offre un buon compromesso tra efficacia nella valutazione della verosimiglianza e
efficienza computazionale, risultando più rapido di Magika nel contesto specifico
considerato, pur perseguendo un obiettivo differente.



\subsection{Variabilità delle Dimensioni dei File Generati}

Poiché la lunghezza dei file influisce direttamente sia sul costo computazionale
del processo di generazione sia sulle rappresentazioni statistiche dei file,
è stata analizzata la distribuzione delle dimensioni dei file generati
dal modello markoviano.

\bigskip

L’obiettivo di questa analisi è verificare empiricamente l’effetto della strategia
di generazione basata su una distribuzione log-uniforme, introdotta nella
Sezione~5.7.4, e valutare se essa produca un insieme di file sufficientemente
eterogeneo ma controllato.

I risultati ottenuti sono riportati nella tabella seguente.

\bigskip

{\centering
\captionof{table}{Statistiche sulla lunghezza dei file generati per ciascun formato.}
\label{tab:file_length_stats}
\begin{tabular}{lcccccc}
\toprule
\textbf{Formato} & \textbf{N} & \textbf{Media} & \textbf{Dev. Std.} & \textbf{Varianza} & \textbf{$L_{\min}$} & \textbf{$L_{\max}$} \\
\midrule
DOCX & 2500 & 8326.5 & 2544.1 & 6\,472\,297.9 & 5305 & 14505 \\
JPG  & 2500 & 4004.6 & 2568.8 & 6\,598\,794.7 & 1024 & 10230 \\
PDF  & 2500 & 4143.9 & 2634.6 & 6\,941\,337.3 & 1028 & 10235 \\
TXT  & 2500 & 3176.6 & 2578.4 & 6\,647\,894.1 & 512  & 10223 \\
\bottomrule
\end{tabular}
\par}



\bigskip

Dall’analisi dei risultati emerge come, per tutti i formati considerati, le
dimensioni dei file generati presentino una variabilità elevata ma controllata.
In particolare, i valori di deviazione standard e varianza risultano comparabili
tra i diversi formati, indicando un comportamento stabile del processo di
generazione.

\bigskip

I valori osservati di $L_{\min}$ e $L_{\max}$ confermano inoltre il rispetto dei
vincoli imposti nella fase di generazione: non vengono prodotti file
eccessivamente piccoli, potenzialmente privi di informazione strutturale, né
file di dimensioni anomale che comporterebbero un costo computazionale
sproporzionato.

\bigskip

Le differenze nelle lunghezze medie tra i formati sono coerenti con gli
intervalli $[L_{\min}, L_{\max}]$ adottati e riflettono le caratteristiche tipiche
dei formati stessi. In particolare, i file DOCX presentano una lunghezza media
maggiore rispetto agli altri formati, in quanto non sono soltanto 
contenitori di testo, ma includono strutture XML e risorse aggiuntive, come spiegato in precedenza.
Mentre i file TXT risultano mediamente più contenuti, come atteso.

\bigskip

Nel complesso, la distribuzione delle lunghezze ottenuta conferma che la scelta
di una distribuzione log-uniforme consente di evitare una concentrazione
artificiale delle dimensioni attorno a un valore fisso, garantendo al contempo
un’adeguata copertura dell’intervallo di generazione.

\bigskip

Analisi ripetute su diverse istanze di generazione, con cardinalità variabile tra
10\,000 e 100\,000 file, hanno mostrato valori statistici analoghi, suggerendo la
robustezza e la riproducibilità del processo di generazione anche rispetto alla
variabilità dimensionale dei file.




\chapter{Conclusioni}

Nel presente lavoro è stato progettato, implementato e validato un sistema per la generazione
di file-esca casuali ma verosimili, finalizzato al supporto di tecniche di Data Flooding against Ransomware.

L'obiettivo principale non è stato la creazione di file semanticamente validi, ma piuttosto la produzione di file
che risultino statisticamente compatibili con i file reali, evitando sia la replica diretta dei dati sensibili,
sia la generazione di file completamente casuali.

\bigskip

A tal fine, è stata introdotta una definizione quantitativa di verosimiglianza statistica, basata su 
rappresentazioni esplicite del contenuto binario e su metriche di distanza e divergenza ampiamente utilizzate
in letteratura.

Su queste basi è stato sviluppato un sistema di rilevamento capace di distinguere tra file reali e file randomici,
mostrando prestazioni elevate in termini di accuratezza e tassi di errore contenuti.


\bigskip

Il sistema di generazione proposto si fonda su modellazioni markoviane del contenuto binario dei file reali,
con un processo di campionamento sviluppato con l'obiettivo di ottimizzare l'efficienza computazionale.

Le scelte progettuali adottate, come la selezione delle porzioni di file da modellare e
l'uso di una distribuzione log-uniforme per la lunghezza dei file generati,
hanno permesso di bilanciare efficacemente la qualità dei file prodotti
con i vincoli computazionali imposti dal contesto operativo.

\bigskip

La valutazione sperimentale ha mostrato che i file generati:
\begin{itemize}
    \item risultano statisticamente distinti dai file completamente randomici;
    \item sono classificati come verosimili dal detector nella maggior parte dei casi;
    \item sono riconosciuti correttamente da strumenti esterni come Magika in una percentuale significativa di casi,
          confermando la coerenza delle firme statistiche apprese, soprattutto per i formati DOCX e TXT.
\end{itemize}

\bigskip

Il confronto prestazionale tra il detector e Magika ha inoltre evidenziato come l'approccio basato
su metriche esplicite e soglie ottimizzate consenta un'elaborazione significativamente più rapida,
rendendo il sistema particolarmente adatto a scenari operativi con grandi quantità di dati.

Questo risultato rafforza la scelta di non adottare modelli di Machine Learning profondi, 
privilegiando invece un approccio più leggero e interpretabile.


\bigskip

In conclusione, il lavoro svolto dimostra come sia possibile manipolare 
in modo controllato l'identità statistica di un file,
generando contenuti artificiali che risultano credibili per analisi automatiche, senza compromettere l’efficienza del sistema. 


\section{Limitazioni e Validità dello Studio}

La nozione di verosimiglianza adottata in questo lavoro è definita nello spazio delle caratteristiche considerate
(ad es. statistiche su n-gram e porzioni strutturali dei file) e mira a riprodurre firme statistiche
compatibili con i campioni reali. Tale criterio non implica correttezza semantica né garantisce piena
conformità agli standard di parsing dei formati. Inoltre, l’efficacia dell’approccio rispetto a
ransomware reali dipende dalle euristiche effettivamente implementate dagli attaccanti e richiede
una validazione in ambiente controllato.

\bigskip

Inoltre, pur utilizzando dataset dell'ordine di alcune migliaia di file per formato, la variabilità osservata potrebbe 
non coprire l'intera eterogeneità presente in scenari reali.

\bigskip


Infine, le prestazioni possono variare tra formati: i risultati sono più solidi per quei formati in cui la struttura osservata presenta pattern ricorrenti, mentre per formati con maggiore complessità binaria e/o componenti ad alta entropia si osserva un margine di distanza residuo.
L'uso di un modello markoviano di ordine 1 (bigrammi) privilegia l'efficienza computazionale, ma può non catturare dipendenze a più lungo raggio, che potrebbero essere rilevanti per alcuni formati.


\section{Sviluppi Futuri}

Una possibile direzione per futuri sviluppi di questo lavoro ricade nello studio di un approccio per il quale 
è possibile convertire un file di uno specifico formato in un altro formato, per riconverterlo successivamente.

\bigskip

Come visto in precedenza, molti ransomware tendono a ignorare file con caratteristiche strutturali atipiche 
usando euristiche basate su analisi statistiche.

In questo contesto, risulta pertanto interessante lo studio degli elementi che definiscono l'identità di un file,
e in particolare in quale misura tali elementi possano essere modificati senza compromettere la riconversione.

\bigskip

Per fare un esempio, si può pensare di avere un file PDF contenente dati sensibili, 
come ad esempio password o chiavi crittografiche.

In questo scenario, l'obiettivo sarebbe quello di convertire il file PDF in un altro formato che sappiamo essere ignorato dai ransomware,
ad esempio un file di log o file altamente entropici.

Questa conversione non ha lo scopo di distruggere il file, ma piuttosto di "mascherarlo", per poi poterlo riconvertire nel formato originale
quando necessario.

\bigskip

Le rappresentazioni basate su N-gram e le metriche di distanza descritte in questo lavoro possono fornire un punto di partenza per questo tipo di analisi.
In particolare, si potrebbe considerare efficace una conversione se la sua distribuzione risulta
statisticamente dissimile da quella del formato originale.

\bigskip

L'esplorazione di questo approccio risulta dunque una strada verso nuove strategie contro le minacce ransomware,
complementari rispetto alle tecniche di Data Flooding e backup tradizionali.

\bigskip

Un ulteriore sviluppo riguarda l’estensione del generatore markoviano a modelli di ordine superiore
(ad es. ordine 2 o 3), al fine di catturare dipendenze a più lungo raggio nel contenuto binario.
Rispetto al modello a bigrammi adottato in questo lavoro, n-grammi più lunghi potrebbero riprodurre
pattern più ricchi e ridurre la distanza residua osservata per formati con maggiore complessità e/o
componenti ad alta entropia. Tale estensione introduce tuttavia un trade-off tra qualità e costo:
l’aumento dell’ordine comporta una crescita significativa dello spazio degli stati, con impatto su memoria
e generazione, oltre a richiedere strategie adatte per gestire la sparsità dei conteggi.
Un’analisi sperimentale sistematica potrebbe quantificare questi effetti confrontando le metriche di verosimiglianza
e i costi runtime, identificando l’ordine minimo in grado di
fornire benefici misurabili nei diversi formati.




\clearpage
\bibliographystyle{IEEEtran} 
\bibliography{reference}






\end{document}
